\chapter{系统辨识 · 鲁棒与频域控制（H∞/μ/灵敏度/回路成形）}
\label{ch:robust-hinf}

% ════════════════════════════════════════════════════════════════
%  控制部 C1a · 第 E 章：辨识 → 鲁棒综合 → 频域多变量验证
%  源：Tutorial 40_控制理论/60_辨识鲁棒与频域 + 150_机器人系统辨识
%  与 ch:control(SISO 频域基础)、C 章(H∞ 状态反馈概览)、D 章(ISS/无源性) 分工
%  记号：θ 为惯性参数(避与圆周率 pi 冲突)、移位算子 z^{-1}、P 为广义被控对象
% ════════════════════════════════════════════════════════════════

\begin{practice}[前置自测]
开始之前，先试答下列问题。若已能流畅作答，可只速读本章的方法论表与速查卡；若卡住，正是本章要补完的。
\begin{enumerate}
  \item 连续时间 LQR 的 Riccati 方程与最优增益 $\mathbf{K}$ 长什么样？$\mathbf{P}$ 矩阵的物理意义是什么？（见 \cref{ch:optimal-control-basics}）
  \item LQG 分离原理说了什么？为什么 Doyle 1978 的反例表明 LQG 的鲁棒裕度可以任意小？
  \item 传递函数的 $H_\infty$ 范数是什么？它在频域的几何含义是什么？
  \item 给定超定线性方程 $\mathbf{A}\mathbf{x}=\mathbf{b}$，最小二乘解的闭式表达是什么？何时唯一？（见 \cref{ch:nlopt}）
  \item 矩阵 $\mathbf{A}$ 的奇异值分解 $\mathbf{A}=\mathbf{U}\boldsymbol\Sigma\mathbf{V}^\top$ 里，$\mathbf{U},\boldsymbol\Sigma,\mathbf{V}$ 各代表什么几何操作？
\end{enumerate}
\end{practice}

\section*{本章目标}
学完本章，你应当能够：
\begin{enumerate}
  \item \textbf{从数据到模型}：用批量/递推最小二乘、预测误差法（PEM）、子空间辨识，把输入输出数据翻译成可控制的模型，并理解持续激励是辨识可行的硬前提；
  \item \textbf{辨识机器人动力学}：把运动方程写成线性回归 $\mathbf{Y}(\mathbf{q},\dot{\mathbf q},\ddot{\mathbf q})\boldsymbol\theta=\boldsymbol\tau$，识别基参数、设计 D-最优激励轨迹、给参数置信区间；
  \item \textbf{数据驱动控制}：理解 Willems 基本引理、DeePC 与 MPC 的等价、System Level Synthesis 的闭环响应参数化；
  \item \textbf{鲁棒综合}：掌握小增益定理、结构奇异值 $\mu$ 与 D--K 迭代、IQC 统一框架，把 $H_\infty$ 从状态反馈推进到输出反馈 DGKF 解与回路成形；
  \item \textbf{频域多变量诊断}：用奇异值方向性、RGA、MIMO 灵敏度、水床积分、广义 Nyquist 分析多变量闭环的耦合、病态方向与基本限制；
  \item \textbf{把三者焊成闭环}：理解“辨识 $\to$ 鲁棒综合 $\to$ 频域验证”为何是可部署最优控制的完整链路。
\end{enumerate}

\begin{note}
\textbf{本章定位与分工.}\;
\cref{ch:optimal-control-basics} 与 \cref{ch:control} 以 LQR/LQG 为中心，建立了\emph{模型已知、噪声为白}的最优控制闭环。本章补完它的三个缺口：模型从哪来（辨识）、不确定下如何保性能（鲁棒综合）、如何验证多变量闭环（频域）。与兄弟章的硬边界：SISO 的灵敏度 $S/T$、增益/相位裕度、Bode 积分魔咒、Nyquist、PID、根轨迹已在 \cref{ch:control} 给出，本章\textbf{一律不重证 SISO 基础}，只做 MIMO 升级与鲁棒综合；$H_\infty$ 的\emph{状态反馈博弈 Riccati 概览}见 \cref{ch:lqr-lqg-riccati}，本章承接并深化到\emph{输出反馈 DGKF 两-Riccati 解 + 混合灵敏度 + $\mu$ 综合 + 回路成形}；ISS、无源性、Port-Hamiltonian 的系统理论归 \cref{ch:lyapunov-stability}，本章的 IQC/KYP 只作“鲁棒判据的频域统一视角”。
\end{note}

\begin{note}
\textbf{本章记号约定盒（三处符号过载，先钉死）.}\;
本章横跨辨识与鲁棒两域，三个符号冲突必须显式消歧：
\begin{enumerate}
  \item \textbf{移位算子 vs 关节角.}\; 辨识传统用 $q$ 记单位移位算子（$qu(t)=u(t+1)$），但机器人动力学里关节角也叫 $q$。本章把移位算子一律写作 $z^{-1}$（后移）或 $z$（前移），关节角保留黑体 $\mathbf{q}$。
  \item \textbf{惯性参数 vs 圆周率.}\; 源用 $\pi$ 记惯性参数向量，但水床积分里 $\pi$ 是圆周率。本章惯性参数一律用 $\boldsymbol\theta$（黑体），圆周率保持细 $\pi$。
  \item \textbf{广义被控对象 vs 物理对象.}\; $\mathbf{P}$ 留给\emph{广义被控对象}（含权函数的 LFT 形式），物理对象记 $\mathbf{G}$（与 \cref{ch:control} 的 $P(s)$ 区别开）。
\end{enumerate}
其余：灵敏度 $\mathbf{S}=(\mathbf{I}+\mathbf{L})^{-1}$、补灵敏度 $\mathbf{T}=\mathbf{L}(\mathbf{I}+\mathbf{L})^{-1}$ 与 \cref{ch:control} 的 \cref{eq:ctrl-sens,eq:ctrl-comp-sens} 同义，MIMO 版为矩阵传递函数，$\mathbf{S}+\mathbf{T}=\mathbf{I}$ 仍成立（\cref{eq:ctrl-S-plus-T}）但 $\mathbf S,\mathbf T$ 不对易；不确定性块 $\boldsymbol\Delta=\mathrm{diag}(\delta_i\mathbf I,\boldsymbol\Delta_j)$ 用标准 $\mu$ 记号。
\end{note}

\noindent\textbf{为什么“辨识 + 鲁棒 + 频域”是最优控制的三块拼图.}\quad
\cref{ch:optimal-control-basics} 的 LQR/LQG 漂亮，但在真实机器人上有三个致命缺口。\emph{缺口一：模型从哪来？}真实机器人不会奉送 $(\mathbf A,\mathbf B,\mathbf C,\mathbf D)$——必须从数据辨识，而辨识偏差直接破坏 LQR 的最优性；一台 18 关节四足，每个关节的摩擦、减速比效率、绕组电阻都未知，CAD 惯量误差常达 $10\%$--$20\%$。\emph{缺口二：LQG 鲁棒裕度可任意小。}Doyle 1978\,\cite{doyle1978guaranteed} 的反例证明：即使 LQG 在名义模型上最优，一个任意小的建模误差就足以让闭环失稳——这不是 bug 而是结构性缺陷，$H_2$ 最优不蕴含鲁棒。回顾 \cref{ch:optimal-control-basics}：LQR 状态反馈有无穷增益裕度与 $60^\circ$ 相位裕度（Kalman 1964\,\cite{kalman1964optimal}），但一旦加观测器变成 LQG，这些保证全部丧失。\emph{缺口三：SISO 工具在 MIMO 失效。}现代机器人是多输入多输出耦合系统，SISO 的 Bode/Nyquist 无法捕捉方向性——pitch 方向稳定性极好的系统可能在 roll 方向几乎失稳。

\begin{insight}
真正可部署的最优控制，必须是\textbf{辨识 $\to$ 鲁棒综合 $\to$ 频域验证}的完整闭环。本章把这三块按统一数学语言（Hankel 矩阵、$H_\infty$ 范数、结构奇异值 $\mu$、闭环响应 $\{\boldsymbol\Phi_x,\boldsymbol\Phi_u\}$）焊接起来：它们其实是同一枚硬币的三面。辨识 = 品尝食材；鲁棒综合 = 食材波动也能出合格菜；频域分析 = 用味觉在各维度上诊断菜品平衡——三者缺一不可。
\end{insight}

% ════════════════════════════════════════════════════════════════
\section{系统辨识范式与最小二乘}
\label{sec:sysid-ls}

\paragraph{动机：CAD 参数不靠谱。}
你面前是一台新机械臂，电机转矩常数 $k_t$、减速器摩擦是 Coulomb 还是 Stribeck、关节柔性刚度 $k_s$——图纸上没有，手册值与实际差 $20\%$。直接拿 CAD 参数设计控制器，在仿真里完美（仿真用的就是 CAD 参数），部署到真机却震荡、跟踪误差大甚至失稳，因为名义模型与真实系统的参数偏差超过了控制器的鲁棒裕度。这正是 sim-to-real gap 的核心来源之一。辨识不是可选步骤，是必要步骤。

\subsection{辨识范式：五步法与四轴分类}
系统辨识把装置的输入输出数据翻译为可用于控制的数学模型，方法论沿四个正交轴划分（\cref{tab:sysid-taxonomy}）。Ljung\,\cite{ljung1999sysid} 的“实验设计 $\to$ 数据采集 $\to$ 模型结构选择 $\to$ 参数估计 $\to$ 模型验证”五步法是业界标准。

\begin{table}[H]\centering\small
\caption{系统辨识的四轴分类}
\label{tab:sysid-taxonomy}
\begin{tabular}{llll}
\toprule
分类轴 & 选项 A & 选项 B & 选择依据 \\
\midrule
参数化 & 参数化模型 $\mathcal M(\boldsymbol\theta)$ & 非参数化（脉冲/频响） & 是否有先验物理结构 \\
域 & 时域（拟合 $y(t)$） & 频域（拟合 $\mathbf G(j\omega)$） & 数据形式与噪声特性 \\
回路 & 开环（$u$ 独立于 $y$） & 闭环（$u$ 由控制器产生） & 系统是否在运行中 \\
知识 & 黑箱（纯数据） & 灰箱（保留已知物理结构） & 物理先验的可靠程度 \\
\bottomrule
\end{tabular}
\end{table}

Ljung 反复强调\textbf{“模型为目的服务”}：仿真、预测与控制三种目的要求不同的辨识准则。最小化一步预测误差的模型未必最适合控制——这正是 Gevers“identification for control”思想的出发点：一个在 one-step 上方差极小、却在 multi-step 上偏差大的模型，不如反之者适合 MPC，因为辨识准则应与控制目标对齐。

\subsection{批量最小二乘：正规方程与几何}
最小二乘是辨识的基石。考虑 ARX 模型（以 $z^{-1}$ 记后移算子）
\begin{equation}
  y(t)+a_1 y(t-1)+\cdots+a_n y(t-n)=b_1 u(t-1)+\cdots+b_m u(t-m)+e(t).
  \label{eq:arx}
\end{equation}
定义回归向量与参数向量
\begin{equation}
  \boldsymbol\varphi(t)=\big[-y(t-1),\dots,-y(t-n),\,u(t-1),\dots,u(t-m)\big]^\top,\quad
  \boldsymbol\theta=[a_1,\dots,a_n,b_1,\dots,b_m]^\top,
\end{equation}
则 \cref{eq:arx} 化为线性回归 $y(t)=\boldsymbol\varphi(t)^\top\boldsymbol\theta+e(t)$。其妙处在于：把动态辨识转化为静态线性回归，因为 $\boldsymbol\varphi(t)$ 只含已观测的过去 $y,u$。堆叠 $N$ 个数据点得超定方程 $\mathbf Y=\boldsymbol\Phi\boldsymbol\theta+\mathbf E$，最小化 $V(\boldsymbol\theta)=\tfrac1N\|\mathbf Y-\boldsymbol\Phi\boldsymbol\theta\|^2$，令梯度为零得正规方程，解得
\begin{equation}
  \hat{\boldsymbol\theta}_N=(\boldsymbol\Phi^\top\boldsymbol\Phi)^{-1}\boldsymbol\Phi^\top\mathbf Y.
  \label{eq:ls-normal}
\end{equation}
解存在且唯一当且仅当 $\boldsymbol\Phi$ 列满秩。物理含义：输入必须“足够丰富”——恒定输入下所有 $\boldsymbol\varphi(t)$ 方向相同，$\boldsymbol\Phi$ 秩为 1，无法分辨 $n+m$ 个参数，这正是 \cref{sec:pe} 持续激励条件的矩阵表述。最小二乘求解的数值实现（QR/LDLT/SVD 而非显式求逆）见 \cref{ch:nlopt}。

\begin{insight}
\cref{eq:ls-normal} 本质是把 $\mathbf Y$ 正交投影到 $\boldsymbol\Phi$ 的列空间。在白噪声假设（$e(t)$ 零均值方差 $\sigma^2$ i.i.d.）下，$\hat{\boldsymbol\theta}_N$ 无偏（$\mathbb E[\hat{\boldsymbol\theta}_N]=\boldsymbol\theta_0$），协方差 $\mathrm{Cov}(\hat{\boldsymbol\theta}_N)=\sigma^2(\boldsymbol\Phi^\top\boldsymbol\Phi)^{-1}$，PE 下一致，且渐近正态 $\sqrt N(\hat{\boldsymbol\theta}_N-\boldsymbol\theta_0)\xrightarrow{d}\mathcal N(0,\sigma^2\mathbf R_\varphi^{-1})$。数据越丰富（$\boldsymbol\Phi^\top\boldsymbol\Phi$ 最小特征值越大），投影越稳定，估计越精确。
\end{insight}

\subsection{递推最小二乘与 Kalman 对偶}
在线辨识中数据逐点到达，每步重算 $(\boldsymbol\Phi^\top\boldsymbol\Phi)^{-1}$ 需 $O(n^3)$，对 1\,kHz 控制循环不可接受。定义逆相关矩阵 $\mathbf P(t)=\big[\sum_{\tau=1}^t\boldsymbol\varphi(\tau)\boldsymbol\varphi(\tau)^\top\big]^{-1}$，注意 $\mathbf P(t)^{-1}=\mathbf P(t-1)^{-1}+\boldsymbol\varphi(t)\boldsymbol\varphi(t)^\top$ 是秩-1 更新，对它用 Sherman--Morrison--Woodbury 引理 $(\mathbf A+\mathbf u\mathbf v^\top)^{-1}=\mathbf A^{-1}-\frac{\mathbf A^{-1}\mathbf u\mathbf v^\top\mathbf A^{-1}}{1+\mathbf v^\top\mathbf A^{-1}\mathbf u}$，得递推最小二乘（RLS）：
\begin{align}
  \mathbf P(t)&=\mathbf P(t-1)-\frac{\mathbf P(t-1)\boldsymbol\varphi(t)\boldsymbol\varphi(t)^\top\mathbf P(t-1)}{1+\boldsymbol\varphi(t)^\top\mathbf P(t-1)\boldsymbol\varphi(t)},\label{eq:rls-P}\\
  \hat{\boldsymbol\theta}(t)&=\hat{\boldsymbol\theta}(t-1)+\mathbf K(t)\big[y(t)-\boldsymbol\varphi(t)^\top\hat{\boldsymbol\theta}(t-1)\big],\quad
  \mathbf K(t)=\frac{\mathbf P(t-1)\boldsymbol\varphi(t)}{1+\boldsymbol\varphi(t)^\top\mathbf P(t-1)\boldsymbol\varphi(t)}.\label{eq:rls-theta}
\end{align}
$\varepsilon(t)=y(t)-\boldsymbol\varphi(t)^\top\hat{\boldsymbol\theta}(t-1)$ 是预测误差，$\mathbf K(t)$ 是修正增益，$\mathbf P(t)$ 是参数协方差。

\begin{theorem}[RLS = 参数估计的 Kalman 滤波]
\label{thm:rls-kalman}
把参数估计写成状态空间：状态方程 $\boldsymbol\theta(t+1)=\boldsymbol\theta(t)$（参数不变，零过程噪声），观测方程 $y(t)=\boldsymbol\varphi(t)^\top\boldsymbol\theta(t)+e(t)$。对此系统应用 Kalman 滤波（\cref{ch:eskf}），所得递推恰为 \cref{eq:rls-P,eq:rls-theta}。即 RLS \emph{就是} Kalman 滤波在参数估计上的特例。
\end{theorem}

\begin{insight}
LQR 与 Kalman 的对偶在辨识中的体现是：\textbf{最优控制（LQR）$\leftrightarrow$ 最优估计（Kalman）$\leftrightarrow$ 最优辨识（RLS）}。三者共享同一个 Riccati/秩-1 更新骨架，只是“被估计的对象”从控制增益换成状态、再换成参数。
\end{insight}

\subsection{遗忘因子与变遗忘因子}
对时变系统旧数据应逐渐被遗忘：四足拾起重物后等效惯量从 $\boldsymbol\theta_{\text{old}}$ 跳到 $\boldsymbol\theta_{\text{new}}$，若不遗忘，旧数据“投票权”太大使估计永远到不了新值。加权损失 $V(\boldsymbol\theta)=\sum_{t=1}^N\lambda^{N-t}\|y(t)-\boldsymbol\varphi(t)^\top\boldsymbol\theta\|^2$，$\lambda\in(0,1)$，最新数据权重 $1$、$k$ 步前权重 $\lambda^k$，等效记忆窗口 $\approx 1/(1-\lambda)$ 步（$\lambda=0.99$ 约 100 步）。带遗忘因子的 RLS 把 \cref{eq:rls-P} 分母的 $1$ 换成 $\lambda$、前乘 $1/\lambda>1$，使 $\mathbf P(t)$ 不再单调趋零，从而保持对新数据的敏感性——这是偏差-方差权衡在辨识中的直接体现。

固定 $\lambda$ 的局限：太小则跟踪快但稳态方差大，太大则反之。\textbf{变遗忘因子}（VFF，Fortescue--Kershenbaum--Ydstie 1981\,\cite{fortescue1981vff}）用预测误差自适应调节
\begin{equation}
  \lambda(t)=1-\frac{\varepsilon(t)^2\,\boldsymbol\varphi(t)^\top\mathbf P(t-1)\boldsymbol\varphi(t)}{c+\boldsymbol\varphi(t)^\top\mathbf P(t-1)\boldsymbol\varphi(t)},
\end{equation}
$|\varepsilon|$ 大时加速遗忘、小时减缓。当参数各分量变化速率不同（摩擦快变、惯量不变），各向同性遗忘低效；\textbf{方向性遗忘}（Kulhavý--Kárný 1984\,\cite{kulhavy1984directional}）只在有新信息的方向上遗忘，避免无激励方向上 $\mathbf P(t)$ 无限增长（协方差 blow-up）。

% ════════════════════════════════════════════════════════════════
\section{预测误差法与模型族}
\label{sec:pem}

\paragraph{从 ARX 到一般噪声模型。}
ARX 简单（线性回归、闭式解），但噪声模型与被控对象耦合：$\mathbf A y=\mathbf B u+e$ 中 $\mathbf A$ 的极点既是系统极点又是噪声极点，若噪声有系统外的动态（如 $60\,$Hz 工频），ARX 会把它当成系统极点产生虚假极点。控制设计需要把“系统在做什么”（确定性动态）与“噪声长什么样”（随机部分）清晰分开，PEM 与更灵活的模型族为此而生。

PEM 的核心是一步预测器。对一般模型 $y(t)=G(z,\boldsymbol\theta)u(t)+H(z,\boldsymbol\theta)e(t)$，一步最优预测为
\begin{equation}
  \hat y(t\mid\boldsymbol\theta)=H^{-1}(z,\boldsymbol\theta)G(z,\boldsymbol\theta)u(t)+\big[1-H^{-1}(z,\boldsymbol\theta)\big]y(t),
  \label{eq:pem-predictor}
\end{equation}
要求 $H$ 最小相位（$H^{-1}$ 稳定）。损失 $V_N(\boldsymbol\theta)=\tfrac1N\sum_t\ell(\varepsilon(t,\boldsymbol\theta))$，$\varepsilon=y-\hat y$，通常取 $\ell=\tfrac12\varepsilon^2$（等价高斯 MLE），$\hat{\boldsymbol\theta}_N=\arg\min V_N$。除 ARX 外这是非线性优化，需 Gauss--Newton。

\begin{theorem}[PEM 渐近高效]
\label{thm:pem-crlb}
在平稳遍历与正则条件下，$\sqrt N(\hat{\boldsymbol\theta}_N-\boldsymbol\theta_0)\xrightarrow{d}\mathcal N(0,\mathbf P_\theta)$，其中 $\mathbf P_\theta=\sigma^2[\mathbb E\,\boldsymbol\psi\boldsymbol\psi^\top]^{-1}$，$\boldsymbol\psi(t,\boldsymbol\theta)=-\partial\varepsilon/\partial\boldsymbol\theta$ 是灵敏度函数。在高斯白噪声下 $\mathbf P_\theta$ 达到 Cramér--Rao 下界，故 PEM 渐近高效（等价 MLE）。
\end{theorem}

四大模型族统一在多项式框架 $A(z)y=\tfrac{B(z)}{F(z)}u+\tfrac{C(z)}{D(z)}e$ 下（\cref{tab:model-family}）。

\begin{table}[H]\centering\small
\caption{四大模型族}
\label{tab:model-family}
\begin{tabular}{lllll}
\toprule
模型族 & 方程 & 动态 $G$ & 噪声 $H$ & 特点 \\
\midrule
ARX & $Ay=Bu+e$ & $B/A$ & $1/A$ & 线性回归、闭式解；$G,H$ 共享 $A$ \\
ARMAX & $Ay=Bu+Ce$ & $B/A$ & $C/A$ & 噪声滤波可调；PEM 非线性 \\
OE & $y=(B/F)u+e$ & $B/F$ & $1$ & 动态与噪声完全解耦；\textbf{控制设计首选} \\
BJ & $y=(B/F)u+(C/D)e$ & $B/F$ & $C/D$ & 最一般；动态/噪声独立参数化 \\
\bottomrule
\end{tabular}
\end{table}

\noindent\textbf{选择指南.}\;只关心控制设计（需精确 $G$）选 OE 或 BJ（动态 $B/F$ 不受噪声模型影响）；还关心预测（需精确 $H$）选 ARMAX 或 BJ；快速原型选 ARX（唯一闭式解，用来定阶）；\emph{闭环辨识必须用 BJ 或精确噪声模型的 ARMAX}，因为此时 $u$ 与噪声相关。

% ════════════════════════════════════════════════════════════════
\section{子空间辨识}
\label{sec:subspace}

\paragraph{绕开非凸迭代。}
PEM 的迭代优化有两个困难：非凸（陷局部极小、初始化敏感）、MIMO 维度暴增。子空间辨识直接从数据的 Hankel 矩阵做 SVD，一次性提取 $(\mathbf A,\mathbf B,\mathbf C,\mathbf D)$，无须初值。

设 $m$ 输入 $p$ 输出，把数据排成块 Hankel 矩阵并按行分过去/未来块 $\mathbf U_p,\mathbf U_f,\mathbf Y_p,\mathbf Y_f$。核心三步：
\begin{enumerate}
  \item \textbf{斜投影}得扩展可观性矩阵：$\mathcal O=\mathbf Y_f/_{\mathbf U_f}\mathbf Z_p$，其中 $\mathbf Z_p=[\mathbf U_p;\mathbf Y_p]$ 是过去信息。几何含义是“把 $\mathbf Y_f$ 投影到过去 $\mathbf Z_p$ 的行空间，但先从 $\mathbf Y_f$ 去除未来输入 $\mathbf U_f$ 的直接影响”——因为 $\mathbf Y_f$ 同时受初始状态（藏在 $\mathbf Z_p$）与未来输入影响，只想提取初始状态贡献。
  \item \textbf{截断 SVD}：$\mathcal O\approx\mathbf U_n\boldsymbol\Sigma_n\mathbf V_n^\top$，阶数由奇异值间隙 $\sigma_n/\sigma_{n+1}$ 判定——前 $n$ 个大（系统模态）、其后突小（噪声）。
  \item \textbf{提取系统矩阵}：扩展可观性 $\boldsymbol\Gamma_i=\mathbf U_n\boldsymbol\Sigma_n^{1/2}$，由 $\mathbf C=\boldsymbol\Gamma_i(1{:}p,:)$、移位性质 $\underline{\boldsymbol\Gamma}_i\mathbf A=\overline{\boldsymbol\Gamma}_i$ 解 $\mathbf A=\underline{\boldsymbol\Gamma}_i^+\overline{\boldsymbol\Gamma}_i$；状态序列 $\hat{\mathbf X}=\boldsymbol\Sigma_n^{1/2}\mathbf V_n^\top$，再对 $\big[\hat{\mathbf x}_{k+1};\mathbf y_k\big]=\big[\begin{smallmatrix}\mathbf A&\mathbf B\\\mathbf C&\mathbf D\end{smallmatrix}\big]\big[\hat{\mathbf x}_k;\mathbf u_k\big]$ 做最小二乘得 $\mathbf B,\mathbf D$。
\end{enumerate}

\begin{table}[H]\centering\small
\caption{三大子空间变体：仅左右加权 $\mathbf W_1\mathcal O\mathbf W_2$ 不同}
\label{tab:subspace}
\begin{tabular}{lll}
\toprule
方法 & 出处 & 特点 \\
\midrule
N4SID & Van Overschee--De Moor 1994\,\cite{vanoverschee1994n4sid} & 斜投影，最通用、数值稳健 \\
MOESP & Verhaegen--Dewilde 1992\,\cite{verhaegen1992moesp} & LQ/RQ 分解 + 工具变量 \\
CVA & Larimore 1990\,\cite{larimore1990cva} & 典型相关加权，高斯下渐近有效 \\
\bottomrule
\end{tabular}
\end{table}

三者唯一区别是 SVD 前乘的左右加权矩阵 $\mathbf W_1,\mathbf W_2$（Van Overschee--De Moor 统一定理）。子空间法与带输入的动态模式分解（DMDc）/Koopman 共骨架：DMDc 等价于子空间辨识的一步版本，Koopman 把非线性系统嵌入提升空间后又回到子空间设置——其正式理论归前沿/RL 部，此处仅作指针。机器人工程师偏爱子空间法，因其非迭代、天然 MIMO、奇异值间隙自动给阶数。

% ════════════════════════════════════════════════════════════════
\section{持续激励条件}
\label{sec:pe}

\paragraph{辨识成功的硬前提。}
给机器人恒定力矩 $u(t)=1$ 辨识不出传递函数——只激励了零频。要辨识 $n$ 阶系统，输入必须在足够多频率上有能量。

\begin{definition}[持续激励的三种等价刻画]
\label{def:pe}
信号 $u$ 是 $n$ 阶持续激励（PE），当且仅当下列之一成立：
\begin{enumerate}
  \item \textbf{Toeplitz 正定}：自相关 Toeplitz 矩阵 $\mathbf R_u^{(n)}=[r_u(i-j)]_{i,j=1}^n\succ0$，其中 $r_u(k)=\lim_{N\to\infty}\tfrac1N\sum_t u(t)u(t-k)$；
  \item \textbf{谱条件}：功率谱 $\Phi_u(\omega)$ 在 $(-\pi,\pi)$ 至少在 $n$ 个不同频点非零；
  \item \textbf{Willems-Hankel}（行为视角）：长度 $T$ 的 $u^d$ 是 $L$ 阶 PE 当且仅当 Hankel 矩阵 $\mathcal H_L(u^d)$ 行满秩（需 $T\ge(m+1)L-1$）。
\end{enumerate}
\end{definition}

一致性定理：辨识 $n$ 阶线性模型需 $u$ 的 PE 阶 $\ge 2n$。常见信号的 PE 阶：白噪声 $\infty$（峰值因子大、可能激励非线性）、PRBS $\ge$ 序列长度、$K$ 频多正弦 $2K$、啁啾近似 $\infty$、常数 $1$、单频正弦 $2$。机器人关节辨识首选多正弦或啁啾，白噪声虽 PE 阶无穷但可能激励限位、不安全。

\begin{insight}
PE 在强化学习中的等价物就是 exploration（$\varepsilon$-greedy / 熵正则 / 好奇心）：都要求输入分布足够“覆盖”以保证参数或价值函数可一致估计。缺 PE 的 RL 陷入 policy collapse（只激励部分模态），与辨识中 $\mathbf R_u$ 奇异、参数不可识别是同一数学现象。
\end{insight}

% ════════════════════════════════════════════════════════════════
\section{机器人动力学辨识}
\label{sec:robot-id}

\paragraph{本章独有重头。}
机器人动力学辨识把 \cref{ch:rigid_body} 的运动方程接到最小二乘与自适应控制上。它的现代理论始于 Atkeson--An--Hollerbach 1986\,\cite{atkeson1986inertial} 的关键发现：虽然 $\mathbf M(\mathbf q)\ddot{\mathbf q}+\mathbf C(\mathbf q,\dot{\mathbf q})\dot{\mathbf q}+\mathbf g(\mathbf q)=\boldsymbol\tau$ 对 $\mathbf q$ 高度非线性，但对惯性参数 $\boldsymbol\theta$（质量、第一矩、惯性矩）\emph{线性}。

\subsection{线性参数化}
每个连杆 $i$ 有 10 个标准惯性参数
\begin{equation}
  \boldsymbol\theta_i=\big[m_i,\,m_i c_{x,i},\,m_i c_{y,i},\,m_i c_{z,i},\,I_{xx,i},I_{xy,i},I_{xz,i},I_{yy,i},I_{yz,i},I_{zz,i}\big]^\top\in\mathbb R^{10},
\end{equation}
注意用\textbf{第一矩} $m_i c_{\bullet,i}$ 而非质心 $c_{\bullet,i}$，正是为了保线性（直接辨识 $c_{x,i}$ 会引入二次项 $m_i c_{x,i}^2$ 破坏线性）。运动方程的每项都可写成“运动学量乘惯性参数之和”：
\begin{equation}
  \boldsymbol\tau=\sum_{i=1}^n\mathbf Y_i(\mathbf q,\dot{\mathbf q},\ddot{\mathbf q})\,\boldsymbol\theta_i=\mathbf Y(\mathbf q,\dot{\mathbf q},\ddot{\mathbf q})\,\boldsymbol\theta,\qquad \mathbf Y\in\mathbb R^{n\times 10n}.
  \label{eq:dyn-regressor}
\end{equation}

\begin{proof}[线性性的证明骨架]
连杆 $i$ 对动能的贡献 $T_i=\tfrac12 m_i\mathbf v_{c_i}^\top\mathbf v_{c_i}+\tfrac12\boldsymbol\omega_i^\top\mathbf I_{c_i}\boldsymbol\omega_i$。由 \cref{ch:rigid_body}，$\mathbf M(\mathbf q)=\sum_i\mathbf J_{v_i}^\top m_i\mathbf J_{v_i}+\mathbf J_{\omega_i}^\top\mathbf R_i\mathbf I_{c_i}\mathbf R_i^\top\mathbf J_{\omega_i}$ 的每个分量是惯性参数的一次齐次函数（雅可比 $\mathbf J_{v_i},\mathbf J_{\omega_i},\mathbf R_i$ 只依赖 $\mathbf q$）。动能 $T=\tfrac12\dot{\mathbf q}^\top\mathbf M(\mathbf q)\dot{\mathbf q}$ 对参数线性，Lagrange 算子 $\tfrac{d}{dt}\tfrac{\partial}{\partial\dot{\mathbf q}}-\tfrac{\partial}{\partial\mathbf q}$ 是线性算子，故广义力对惯性参数仍线性。
\end{proof}

线性参数化把辨识从非凸优化（多局部极值、初值敏感、计算量大）变成凸最小二乘（全局最优、闭式解、$O(n^3)$）。

\subsection{基参数}
$10n$ 个标准参数中很多\textbf{不可辨识}，两类原因：(A) \emph{线性依赖}——某些组合始终以同一线性组合出现（如 $I_1+m_1 l_{c1}^2$，无法从力矩分开）；(B) \emph{对力矩无贡献}——如固定于基座的首连杆纯质量 $m_1$ 乘零力臂。对完整回归矩阵 $\mathbf Y_{\text{full}}$（多采样点堆叠）做列空间分析，$\mathrm{rank}(\mathbf Y_{\text{full}})=p\le 10n$，存在 $p$ 个独立组合即\textbf{基参数} $\boldsymbol\theta_b\in\mathbb R^p$ 使 $\mathbf Y_{\text{full}}\boldsymbol\theta=\mathbf Y_b\boldsymbol\theta_b$ 且 $\mathbf Y_b$ 满秩。

\begin{algo}[Gautier 基参数提取]
\label{alg:base-param}
在多个随机构型处算 $\mathbf Y_k$ 堆叠成 $\bar{\mathbf Y}$；对 $\bar{\mathbf Y}$ 做列主元 QR 分解 $\bar{\mathbf Y}\mathbf P=\mathbf Q\mathbf R$；$\mathbf R$ 前 $p$ 个对角元 $|R_{ii}|>\epsilon$ 对应基参数“骨架”，剩余 $10n-p$ 个或为零或可吸收进基参数。\emph{用 QR 列主元而非 SVD}：SVD 给正交基会混合不同连杆的参数，破坏物理可解释性（Gautier--Khalil\,\cite{gautier1991identifiable}）。
\end{algo}

典型规模随拓扑而变、以数值秩为准：7-DOF 机械臂 70 个标准参数，基参数通常在 40--60 之间（加摩擦后更多，随关节类型与结构变化）。摩擦的 Coulomb+viscous 模型 $\tau_{\text{fric},i}=f_{v,i}\dot q_i+f_{c,i}\,\mathrm{sign}(\dot q_i)$ 也线性、扩展回归矩阵；零速附近用 $\tanh(\alpha\dot q_i)$（$\alpha\sim50$--$200$）光滑 $\mathrm{sign}$ 以避不连续，且 $\tanh$ 只依赖 $\dot q$ 不依赖参数，仍保线性。Stribeck 模型对 $v_s$ 非线性，需先辨 Coulomb+viscous 再单拟 Stribeck。

\begin{pitfall}[实现注意：URDF 惯性须转到连杆原点]
URDF 的 \texttt{inertia} 标签是关于\emph{质心}的惯性张量，回归矩阵用的是关于\emph{连杆原点}的。二者由平行轴定理关联 $\mathbf I_{\text{origin}}=\mathbf I_{\text{com}}+m(\mathbf c^\top\mathbf c\,\mathbf I_3-\mathbf c\mathbf c^\top)$。漏掉转换会使辨识残差不收敛、参数偏差极大且 bug 隐蔽。此外 $\ddot{\mathbf q}$ 几乎无传感器直接测量，应由 Fourier 轨迹解析导数给出（数值微分放大噪声，见 \cref{sec:freq-id}）。
\end{pitfall}

\subsection{OLS/WLS/IRLS 与 Gauss--Markov}
$K$ 个时刻堆叠得超定方程 $\bar{\mathbf Y}\boldsymbol\theta_b=\bar{\boldsymbol\tau}+\bar{\boldsymbol\epsilon}$。OLS 解 $\hat{\boldsymbol\theta}_b=(\bar{\mathbf Y}^\top\bar{\mathbf Y})^{-1}\bar{\mathbf Y}^\top\bar{\boldsymbol\tau}$，Hessian $2\bar{\mathbf Y}^\top\bar{\mathbf Y}\succ0$ 保证唯一全局最小。若噪声异方差（不同关节力矩传感器精度不同），用 WLS：$\hat{\boldsymbol\theta}_b^{\text{WLS}}=(\bar{\mathbf Y}^\top\boldsymbol\Sigma^{-1}\bar{\mathbf Y})^{-1}\bar{\mathbf Y}^\top\boldsymbol\Sigma^{-1}\bar{\boldsymbol\tau}$；$\boldsymbol\Sigma$ 由 OLS 残差迭代估计即 IRLS。

\begin{theorem}[Gauss--Markov：OLS 为 BLUE]
\label{thm:gauss-markov}
线性回归 $\bar{\boldsymbol\tau}=\bar{\mathbf Y}\boldsymbol\theta_b+\bar{\boldsymbol\epsilon}$ 中若 $\mathbb E[\bar{\boldsymbol\epsilon}]=0$ 且 $\mathrm{Cov}(\bar{\boldsymbol\epsilon})=\sigma^2\mathbf I$（同方差不相关），则 OLS 估计在所有线性无偏估计中方差最小（BLUE）。
\end{theorem}
\begin{proof}[骨架]
设任意线性无偏估计 $\tilde{\boldsymbol\theta}_b=\mathbf A\bar{\boldsymbol\tau}$，无偏即 $\mathbf A\bar{\mathbf Y}=\mathbf I$。令 $\mathbf D=\mathbf A-(\bar{\mathbf Y}^\top\bar{\mathbf Y})^{-1}\bar{\mathbf Y}^\top$，由 $\mathbf D\bar{\mathbf Y}=0$ 交叉项消去，得 $\mathrm{Cov}(\tilde{\boldsymbol\theta}_b)=\sigma^2(\bar{\mathbf Y}^\top\bar{\mathbf Y})^{-1}+\sigma^2\mathbf D\mathbf D^\top\succeq\mathrm{Cov}(\hat{\boldsymbol\theta}_b)$，等号当且仅当 $\mathbf D=0$。
\end{proof}

进一步，对高斯噪声 Fisher 信息 $\mathcal I(\boldsymbol\theta_b)=\sigma^{-2}\bar{\mathbf Y}^\top\bar{\mathbf Y}$，Cramér--Rao 下界 $\mathrm{Cov}(\hat{\boldsymbol\theta}_b)\succeq\sigma^2(\bar{\mathbf Y}^\top\bar{\mathbf Y})^{-1}$ 恰被 OLS 取到——高斯下 OLS 渐近有效（等价 MLE）。

\subsection{激励轨迹设计}
辨识精度不仅靠算法更靠实验设计。$\hat{\boldsymbol\theta}_b$ 的协方差正比于 $(\bar{\mathbf Y}^\top\bar{\mathbf Y})^{-1}$，故要最大化 $\bar{\mathbf Y}^\top\bar{\mathbf Y}$ 的“大小”。三种最优性准则（\cref{tab:doptimal}）。

\begin{table}[H]\centering\small
\caption{激励轨迹最优性准则}
\label{tab:doptimal}
\begin{tabular}{llll}
\toprule
准则 & 目标 & 几何含义 & 适用 \\
\midrule
D-最优 & $\max\det(\bar{\mathbf Y}^\top\bar{\mathbf Y})$ & 最小化协方差椭球体积 & 所有参数同等重要 \\
A-最优 & $\min\mathrm{tr}((\bar{\mathbf Y}^\top\bar{\mathbf Y})^{-1})$ & 最小化方差之和 & 关注总体精度 \\
E-最优 & $\max\lambda_{\min}(\bar{\mathbf Y}^\top\bar{\mathbf Y})$ & 最小化最大方差 & 关注最差参数 \\
\bottomrule
\end{tabular}
\end{table}

实践最常用 D-最优。Swevers 1997\,\cite{swevers1997optimal} 把轨迹参数化为有限阶 Fourier 级数
\begin{equation}
  q_j(t)=q_{0,j}+\sum_{l=1}^{N_h}\Big(\frac{a_{l,j}}{l\omega_f}\sin(l\omega_f t)-\frac{b_{l,j}}{l\omega_f}\cos(l\omega_f t)\Big),
  \label{eq:fourier-traj}
\end{equation}
$\omega_f=2\pi/T_f$ 是基频。除以 $l\omega_f$ 是为让导数简洁：$\dot q_j=\sum_l(a_{l,j}\cos+b_{l,j}\sin)$，$\ddot q_j=\sum_l l\omega_f(-a_{l,j}\sin+b_{l,j}\cos)$——\emph{速度的 Fourier 系数就是 $a,b$}（优化速度幅度即优化系数），加速度幅度正比于 $l\omega_f$（高次谐波力矩更大），且轨迹自动周期化（$\ddot{\mathbf q}$ 由解析公式给出，避数值微分）。优化问题
\begin{equation}
  \max_{\boldsymbol\xi}\ \log\det\big(\bar{\mathbf Y}(\boldsymbol\xi)^\top\bar{\mathbf Y}(\boldsymbol\xi)\big)
  \quad\text{s.t.}\quad q_{\min}\le q(t_k;\boldsymbol\xi)\le q_{\max},\ |\dot q(t_k;\boldsymbol\xi)|\le\dot q_{\max},
\end{equation}
用 $\log\det$ 避数值溢出（$\det$ 可达 $10^{100}$）且为凹（整体仍非凸）；这是非线性约束优化，用 SQP 或 CMA-ES 求解（多随机初始化，取条件数 $<100$ 者即可，不必全局最优）。

\begin{insight}
辨识精度应匹配控制需求，边际收益递减：参数误差从 $50\%$ 降到 $5\%$ 改善巨大，从 $5\%$ 到 $1\%$ 可能几乎无感。MPC/前馈对参数敏感（需 $<5\%$），PD+重力补偿不敏感（$<20\%$ 即可）。关键是信息矩阵 $\bar{\mathbf Y}^\top\bar{\mathbf Y}$ 的条件数而非数据量——100 万共线点不如 100 个正交点。
\end{insight}

\noindent\textbf{前向钩子.}\;本节的回归矩阵 $\mathbf Y$ 在自适应控制（Slotine--Li 1987\,\cite{slotineli1987adaptive}，律 $\dot{\hat{\boldsymbol\theta}}=-\boldsymbol\Gamma\mathbf Y^\top\mathbf s$）、操作空间阻抗（\cref{ch:operational-space}）与 WBC/TSID（\cref{ch:force-wbc}）中\emph{复用同一个 $\mathbf Y$}——辨识质量直接决定自适应控制的收敛速度与稳态精度。

% ════════════════════════════════════════════════════════════════
\section{频域辨识与柔性关节}
\label{sec:freq-id}

\paragraph{频段质量不均，应选择性利用。}
时域方法把所有频率混在一起，但低频（$<1\,$Hz，重力/Coulomb 摩擦主导）信噪比高、中频（1--10\,Hz，惯性/Coriolis 显著）是“黄金频段”、高频（$>10\,$Hz，噪声/柔性振动主导）应避开。频域辨识允许选择性加权特定频段。

多正弦激励 $u_j(t)=\sum_{l=1}^{N_h}A_{l,j}\cos(2\pi f_l t+\phi_{l,j})$ 在频域只在选定频点有能量。\textbf{Schroeder 相位} $\phi_{l,j}=\phi_{1,j}-l(l-1)\pi/N_h$ 使波峰因子接近 $\sqrt2$（正弦理论下限），故在力矩限制下可用更大激励幅度。用\emph{奇数次谐波} $f_l=(2l-1)\Delta f$ 还能检测非线性：线性系统对奇数谐波输入只在奇数谐波响应，偶数谐波出现能量即暴露非线性。频率响应估计用 ETFE $\hat{\mathbf G}(j\omega_k)=\mathbf Y(\omega_k)/\mathbf U(\omega_k)$（方差不随数据长度降，需 Welch 分段平均，谱估计见 \cref{ch:eskf}），再拟合物理模型参数。

\subsection{柔性关节双惯量模型}
高性能关节非刚性，谐波减速器引入电机侧与负载侧的弹性耦合，使传递函数从积分器变为含共振-反共振对（Spong 1987\,\cite{spong1987flexible}）。双惯量模型
\begin{equation}
  J_m\ddot\theta_m+k_s(\theta_m-\theta_l)=\tau_m,\qquad J_l\ddot\theta_l-k_s(\theta_m-\theta_l)=-\tau_{\text{ext}},
\end{equation}
从电机力矩 $\tau_m$ 到负载角 $\theta_l$ 的传递函数及其特征频率为
\begin{equation}
  G(s)=\frac{\theta_l}{\tau_m}=\frac{k_s}{J_m J_l s^4+k_s(J_m+J_l)s^2},\quad
  \omega_a=\sqrt{\frac{k_s}{J_l}},\quad
  \omega_r=\sqrt{\frac{k_s(J_m+J_l)}{J_m J_l}}.
  \label{eq:flex-joint}
\end{equation}
$\omega_a$ 是负载侧看到的反共振（增益突降），$\omega_r$ 是共振（增益突升）。谐波减速器低阻尼（$\zeta\sim0.01$--$0.05$）+ 中等共振频率（30--100\,Hz）使控制器带宽受严重限制：若 PD 增益交叉频率接近 $\omega_r$ 即共振震荡。大振幅下刚度非线性（共振频率随振幅变），辨识须在工作点附近用小振幅激励。\cref{fig:flex-bode} 示意其奇异值/Bode 形状。

\begin{figure}[ht]
\centering
\begin{tikzpicture}[>=Stealth, scale=1.0]
  % axes
  \draw[->, thick] (0,0) -- (8.4,0) node[right]{$\log\omega$};
  \draw[->, thick] (0,-2.0) -- (0,2.4) node[above]{$|G(j\omega)|$ (dB)};
  % rigid -40dB/dec baseline (double integrator) then notch + resonance
  \draw[main, very thick]
    (0.4,1.9) -- (2.6,0.2)               % -40 dB/dec slope
    .. controls (3.0,-0.4) and (3.1,-1.5) .. (3.3,-1.6) % anti-resonance dip
    .. controls (3.55,-1.5) and (3.6,1.4) .. (3.75,1.5) % resonance peak
    .. controls (4.0,0.8) and (4.4,-0.3) .. (7.8,-1.7); % roll off
  % markers
  \draw[dashed, gray] (3.3,-1.6) -- (3.3,-2.0) node[below, black, font=\small]{$\omega_a$};
  \draw[dashed, gray] (3.75,1.5) -- (3.75,-2.0) node[below, black, font=\small]{$\omega_r$};
  \node[main, font=\small] at (1.4,1.4) {$-40$ dB/dec};
  \node[font=\small, align=center] at (3.05,-1.95) {};
  \node[font=\small] at (2.55,-1.3) {反共振};
  \node[font=\small] at (4.55,1.6) {共振};
\end{tikzpicture}
\caption{柔性关节双惯量模型的幅频示意：刚体段 $-40\,$dB/dec，$\omega_a$ 处反共振凹陷，$\omega_r$ 处共振尖峰。控制带宽须避开 $\omega_r$。}
\label{fig:flex-bode}
\end{figure}

\subsection{Notch 滤波与 H∞ 替代}
辨出 $\omega_r$ 后标准做法是加 notch（陷波器）
\begin{equation}
  N(s)=\frac{s^2+2\zeta_z\omega_n s+\omega_n^2}{s^2+2\zeta_p\omega_n s+\omega_n^2},\qquad \omega_n=\omega_r,
\end{equation}
$\zeta_z$ 小则 notch 深、$\zeta_p$ 大则 notch 宽。\textbf{带宽-鲁棒权衡}：notch 太窄对频率偏移敏感（$\omega_r$ 温漂或辨识误差会“错过”共振），太宽则附近相位损失大降低稳定裕度。离散化对窄带极敏感，须用\emph{预畸变 Tustin}：以 $\omega_d=(2/T_s)\tan(\omega_n T_s/2)$ 替代 $\omega_n$ 设计，保证中心频率离散化后不偏移。共振频率随工况变时用自适应 notch（查表或窄带 RLS 在线估计 $\omega_r$）。

更系统的替代是把柔性当\emph{不确定性}，用 $H_\infty$ 混合灵敏度（见 \cref{sec:hinf-synth}）让优化器自动避开共振：权函数 $W_3(s)$ 在 $\omega_r$ 附近设大值，告诉优化器“此频段不确定性大、勿用高增益”。这与单纯加 notch 相比的好处是把多对共振、相位损失与鲁棒裕度统一进一个优化。柔性关节的 notch 与 $H_\infty$ 综合在力控部柔性关节控制中续用（\cref{ch:force-wbc}）。

% ════════════════════════════════════════════════════════════════
\section{数据驱动控制：Willems、DeePC 与 SLS}
\label{sec:data-driven}

\paragraph{跳过辨识，直接从数据控制。}
传统范式是“先辨识 $(\mathbf A,\mathbf B,\mathbf C,\mathbf D)$ 再设计控制器”，但辨识本身有误差。Willems 基本引理给出另一条路：\emph{一段足够丰富的数据本身就是系统的非参数化表示}。

\begin{theorem}[Willems 基本引理]
\label{thm:willems}
设可控 LTI 系统阶为 $n$、输入维度 $m$。若离线数据输入 $u^d$ 是 $(L+n)$ 阶 PE，则系统所有长度 $L$ 的输入输出轨迹恰为 Hankel 矩阵的列张成空间（Willems--Rapisarda--Markovsky--De Moor 2005\,\cite{willems2005persistency}）：
\begin{equation}
  \mathrm{colspan}\begin{bmatrix}\mathcal H_L(u^d)\\\mathcal H_L(y^d)\end{bmatrix}=\mathfrak B_L,
\end{equation}
其中 $\mathfrak B_L$ 是系统行为集的长度-$L$ 截段。
\end{theorem}

核心含义：任何合法未来轨迹 $(u,y)$ 可表示为数据 Hankel 各列的线性组合——\emph{数据本身就是模型}。\textbf{DeePC}（Coulson--Lygeros--Dörfler，ECC 2019\,\cite{coulson2019deepc}）把它嵌入 MPC，将 Hankel 切为过去/未来块：
\begin{equation}
  \min_{\mathbf g,\mathbf u,\mathbf y}\ \sum_k\|\mathbf y_k-\mathbf r_k\|_{\mathbf Q}^2+\|\mathbf u_k\|_{\mathbf R}^2
  \ \ \text{s.t.}\ \
  \begin{bmatrix}\mathbf U_p\\\mathbf Y_p\\\mathbf U_f\\\mathbf Y_f\end{bmatrix}\mathbf g
  =\begin{bmatrix}\mathbf u_{\text{ini}}\\\mathbf y_{\text{ini}}\\\mathbf u\\\mathbf y\end{bmatrix},\quad
  \mathbf u\in\mathcal U,\ \mathbf y\in\mathcal Y.
  \label{eq:deepc}
\end{equation}

\begin{theorem}[DeePC 与 MPC 等价]
\label{thm:deepc-mpc}
对确定性 LTI 系统，设 $T_{\text{ini}}\ge n$、数据满足 PE，则 \cref{eq:deepc} 与基于真实模型的 MPC 给出相同的最优控制序列。
\end{theorem}
\begin{proof}[骨架]
由 \cref{thm:willems} 任何合法轨迹可表为 $[\mathbf U;\mathbf Y]\mathbf g$；$T_{\text{ini}}\ge n$ 保证初始条件约束 $[\mathbf U_p;\mathbf Y_p]\mathbf g=[\mathbf u_{\text{ini}};\mathbf y_{\text{ini}}]$ 与给定 $\mathbf x_0$ 等价（过去 I/O 可定 $\mathbf x_0$）；未来约束等价于沿 $(\mathbf A,\mathbf B,\mathbf C,\mathbf D)$ 从 $\mathbf x_0$ 演化。故两问题可行集相同、目标相同、解相同。
\end{proof}

\begin{pitfall}[噪声下必须正则化]
确定性等价仅在 LTI 无噪时成立。噪声使真实初始条件不在数据 Hankel 列空间中，精确等式约束\emph{无可行解}。加松弛 $\sigma_y$ 与正则化（Coulson--Lygeros--Dörfler，CDC 2019\,\cite{coulson2019regularized}）$+\lambda_g\|\mathbf g\|^2+\lambda_y\|\boldsymbol\sigma_y\|^2$ 后变近似匹配；$\lambda_g$ 太大退化为“选最近数据”、太小过拟合噪声，经验法则 $\lambda_g\sim\sigma_e^2$、$\lambda_y\sim10\sigma_e^2$。在 Wasserstein 分布鲁棒优化（DRO）视角下，正则化等价于对最坏情况不确定性的保护。DeePC$\leftrightarrow$Tube/鲁棒 MPC 的展开见 \cref{ch:mpc-advanced}。
\end{pitfall}

\subsection{System Level Synthesis}
传统控制设计的自由度是控制器 $\mathbf K$，但 $\mathbf K$ 通过反馈非线性地影响闭环传递函数，使稀疏/局部/FIR 等约束无法凸施加。SLS（Anderson--Doyle--Low--Matni\,\cite{anderson2019sls}）把自由度换成\textbf{闭环响应}本身。对 $\mathbf x_{k+1}=\mathbf A\mathbf x_k+\mathbf B\mathbf u_k+\mathbf w_k$，闭环满足 $[\mathbf x;\mathbf u]=[\boldsymbol\Phi_x;\boldsymbol\Phi_u]\mathbf w$，其中 $\boldsymbol\Phi_x=(z\mathbf I-\mathbf A-\mathbf B\mathbf K)^{-1}$、$\boldsymbol\Phi_u=\mathbf K(z\mathbf I-\mathbf A-\mathbf B\mathbf K)^{-1}$。实现性约束是仿射（凸）的：
\begin{equation}
  [z\mathbf I-\mathbf A,\ -\mathbf B]\begin{bmatrix}\boldsymbol\Phi_x\\\boldsymbol\Phi_u\end{bmatrix}=\mathbf I,
  \qquad \boldsymbol\Phi_x,\boldsymbol\Phi_u\in\tfrac1z\mathcal{RH}_\infty,
  \label{eq:sls-affine}
\end{equation}
它参数化\emph{所有内稳定闭环响应}，控制器由 $\mathbf K=\boldsymbol\Phi_u\boldsymbol\Phi_x^{-1}$ 恢复。SLS 推广了 Youla 参数化（经典 Youla $\mathbf K=(\mathbf Y-\mathbf M\mathbf Q)(\mathbf X-\mathbf N\mathbf Q)^{-1}$ 需已知一个初始稳定控制器、约束在 $\mathbf Q$ 空间可能非凸），而 SLS 直接从 $\mathbf A,\mathbf B$ 出发、约束始终凸、天然支持系统级约束（FIR 截断、稀疏 $\boldsymbol\Phi_u$ 对应分布式控制——四足每条腿只见本腿传感器，稀疏带状约束 = 每条腿只影响邻近腿）。

\begin{insight}
SLS 给出学习-控制的有限样本保证（Dean--Mania--Matni--Recht--Tu 2020\,\cite{dean2020sample}）：从 $N$ 个样本辨识 $(\hat{\mathbf A},\hat{\mathbf B})$ 得误差集，在其上做鲁棒综合，给出 LQR 相对代价的有限样本\emph{次优界}（误差量级 $\sim O(1/\sqrt N)$）。这被视为首个端到端有限样本 LQR 次优保证之一——但应给次优界口径，勿绝对化。
\end{insight}

% ════════════════════════════════════════════════════════════════
\section{鲁棒控制 I：小增益、μ 与 IQC}
\label{sec:robust-mu}

\paragraph{从“模型已知”到“不确定下保性能”。}
辨识给出名义模型，但总有残差。鲁棒控制的任务是：在一族模型（不确定集）上保证稳定与性能。统一工具是线性分式变换（LFT）：把不确定性 $\boldsymbol\Delta$ 与控制器 $\mathbf K$ 从对象中“拉出来”。

\subsection{LFT 建模}
设广义被控对象 $\mathbf M$ 分块 $\mathbf M=\big[\begin{smallmatrix}\mathbf M_{11}&\mathbf M_{12}\\\mathbf M_{21}&\mathbf M_{22}\end{smallmatrix}\big]$。\textbf{下 LFT}（接控制器/反馈在下通道）与\textbf{上 LFT}（接不确定性在上通道）：
\begin{align}
  \mathcal F_l(\mathbf M,\mathbf K)&=\mathbf M_{11}+\mathbf M_{12}\mathbf K(\mathbf I-\mathbf M_{22}\mathbf K)^{-1}\mathbf M_{21},\label{eq:lower-lft}\\
  \mathcal F_u(\mathbf M,\boldsymbol\Delta)&=\mathbf M_{22}+\mathbf M_{21}\boldsymbol\Delta(\mathbf I-\mathbf M_{11}\boldsymbol\Delta)^{-1}\mathbf M_{12}.\label{eq:upper-lft}
\end{align}
\textbf{本章统一约定}（与 Zhou--Doyle--Glover\,\cite{zhou1996robust} 标准一致）：性能/不确定通道在上、控制通道在下时，下 LFT 接 $\mathbf K$、上 LFT 接 $\boldsymbol\Delta$。整个鲁棒分析归结为 $\boldsymbol\Delta$--$\mathbf M$--$\mathbf K$ 结构：先用 $\mathbf K$ 闭下通道得 $\mathbf M(s)=\mathcal F_l(\mathbf P,\mathbf K)$，再分析 $\mathbf M$ 与 $\boldsymbol\Delta$ 的上 LFT 互联稳定性。

\subsection{小增益定理}
\begin{theorem}[小增益定理]
\label{thm:small-gain}
若 $\mathbf M,\boldsymbol\Delta\in\mathcal{RH}_\infty$ 且 $\|\mathbf M\|_\infty\cdot\|\boldsymbol\Delta\|_\infty<1$，则反馈回路 $(\mathbf I-\mathbf M\boldsymbol\Delta)^{-1}$ 内稳定（Zames 1966\,\cite{zames1966inputoutput}；Desoer--Vidyasagar 1975\,\cite{desoer1975feedback}）。
\end{theorem}
$\|\mathbf M\|_\infty$ 是 $\mathbf M$ 放大信号的最大倍数，乘积小于 1 即“信号绕一圈回来变小”，故稳定；鲁棒稳定裕度为 $1/\|\mathbf M\|_\infty$。这像音响啸叫判据：麦克风拾取扬声器声（增益 $\mathbf M$）经房间反射（$\boldsymbol\Delta$）回到麦克风，$\mathbf M\boldsymbol\Delta<1$ 则每循环衰减、不啸叫。

\textbf{当 $\boldsymbol\Delta$ 有结构时小增益过于保守}：满块复不确定性下小增益充要，但实际 $\boldsymbol\Delta$ 常是块对角 $\mathrm{diag}(\delta_1\mathbf I,\delta_2\mathbf I,\boldsymbol\Delta_3)$，把它当满块处理忽略了结构信息——这正是引入 $\mu$ 的动机。

\subsection{结构奇异值 μ}
\begin{definition}[结构奇异值]
\label{def:mu}
\begin{equation}
  \mu_{\boldsymbol\Delta}(\mathbf M)=\frac{1}{\min\{\bar\sigma(\boldsymbol\Delta):\det(\mathbf I-\mathbf M\boldsymbol\Delta)=0,\ \boldsymbol\Delta\in\boldsymbol\Delta\}},
\end{equation}
若无此 $\boldsymbol\Delta$ 则 $\mu=0$（Doyle 1982\,\cite{doyle1982mu}）。$\mu(\mathbf M)$ 是使闭环失稳所需最小结构化不确定性大小的倒数，$\mu$ 越大系统越脆弱。
\end{definition}

\begin{theorem}[μ 鲁棒稳定判据]
\label{thm:mu-rs}
闭环对结构化 $\boldsymbol\Delta$ 鲁棒稳定 $\iff \sup_\omega\mu_{\boldsymbol\Delta}(\mathbf M(j\omega))<1$。
\end{theorem}

$\mu$ 精确计算 NP-hard，实用上夹逼：\textbf{下界}用功率算法；\textbf{D-scaling 上界}
\begin{equation}
  \mu(\mathbf M)\le\inf_{\mathbf D\in\mathcal D}\bar\sigma(\mathbf D\mathbf M\mathbf D^{-1}),
  \label{eq:d-scaling}
\end{equation}
其中 $\mathbf D$ 与 $\boldsymbol\Delta$ 结构对易。该上界的精确性结论（块数 $2K+L\le3$ 时上界 $=\mu$）归 \textbf{Packard--Doyle 1993}\,\cite{packard1993mu}（注意：$\mu$ 的\emph{定义}归 Doyle 1982，D-scaling 上界理论归 Packard--Doyle 1993，二者勿混）。更多块可能存在 gap，上下界接近时可信任上界。\cref{fig:mu-vs-sigma} 示意 $\mu$ 如何位于 $\underline\sigma$ 与 $\bar\sigma$ 之间。

\begin{figure}[ht]
\centering
\begin{tikzpicture}[>=Stealth, scale=1.0]
  \draw[->, thick] (0,0) -- (8.4,0) node[right]{$\omega$};
  \draw[->, thick] (0,0) -- (0,3.2) node[above]{增益};
  % upper bound bar-sigma
  \draw[main, very thick] (0.3,2.6) .. controls (3,2.9) and (5,1.6) .. (8,1.0)
     node[right, main]{$\bar\sigma(\mathbf M)$ (满块小增益)};
  % mu curve in between
  \draw[second, very thick, dashed] (0.3,1.7) .. controls (3,2.0) and (5,1.05) .. (8,0.62)
     node[right, second]{$\mu_{\boldsymbol\Delta}(\mathbf M)$ (结构)};
  % lower bound under-sigma
  \draw[third, very thick, dotted] (0.3,0.9) .. controls (3,1.1) and (5,0.55) .. (8,0.32)
     node[right, third]{$\underline\sigma$ 下界};
  \draw[dashed,gray] (0,1.0) -- (8.0,1.0);
  \node[font=\small] at (1.0,1.15) {$\mu{=}1$};
\end{tikzpicture}
\caption{$\underline\sigma\le\mu_{\boldsymbol\Delta}\le\bar\sigma$：满块小增益用 $\bar\sigma$ 保守，$\mu$ 利用结构信息更紧。鲁棒稳定要求 $\sup_\omega\mu<1$。}
\label{fig:mu-vs-sigma}
\end{figure}

\textbf{鲁棒性能}：鲁棒稳定只保证闭环不崩、性能可能任意差。鲁棒性能等价于把性能不确定性 $\boldsymbol\Delta_P$（虚拟满块）加入结构：$\sup_\omega\mu_{\boldsymbol\Delta\cup\boldsymbol\Delta_P}(\mathbf M(j\omega))<1$。

\subsection{D--K 迭代 μ 综合}
\begin{algo}[D--K 迭代]
\label{alg:dk}
交替优化求 $\mu$ 综合控制器：\textbf{K-step} 固定 $\mathbf D(\omega)$，运行 $H_\infty$ 综合解 $\min_{\mathbf K}\|\mathbf D\,\mathcal F_l(\mathbf P,\mathbf K)\,\mathbf D^{-1}\|_\infty$；\textbf{D-step} 固定 $\mathbf K$，逐频率解 $\min_{\mathbf D}\bar\sigma(\mathbf D\mathbf M\mathbf D^{-1})$（凸）；\textbf{拟合} 用稳定最小相位有理 $\mathbf D(s)$ 拟合频率数据；迭代至收敛。\emph{不保证全局最优}（交替优化、可能局部最优），实践从不同初始 $\mathbf D$ 多次运行。
\end{algo}

利用结构信息使 $\mu$ 综合比直接 $H_\infty$ 更少保守：对末端负载不确定的 2-DOF 机械臂，PID 在最大负载下失稳、LQR 大超调、$H_\infty$（忽略结构）超调 $25\%$、$\mu$ 综合（用结构）超调 $15\%$。

\subsection{IQC 框架}
小增益与 $\mu$ 限于线性不确定性，但真实机器人有大量非线性不确定性（饱和、死区、时变延迟、神经网络）。积分二次约束（IQC）提供最一般框架。

\begin{theorem}[IQC 主定理]
\label{thm:iqc}
算子 $\boldsymbol\Delta$ 满足由 $\boldsymbol\Pi$ 定义的 IQC，若对所有 $(v,w=\boldsymbol\Delta v)$ 有 $\int_{-\infty}^\infty[\hat v;\hat w]^*\boldsymbol\Pi(j\omega)[\hat v;\hat w]\,d\omega\ge0$。若反馈回路 $(\mathbf G,\boldsymbol\Delta)$ 满足：(1) 对所有 $\tau\in[0,1]$，$\tau\boldsymbol\Delta$ 的回路良定（同伦条件）；(2) 存在 $\boldsymbol\Pi$ 使 $\boldsymbol\Delta$ 满足 IQC；(3) 频域不等式 $[\mathbf G(j\omega);\mathbf I]^*\boldsymbol\Pi(j\omega)[\mathbf G(j\omega);\mathbf I]\prec0\ \forall\omega$，则闭环内稳定（Megretski--Rantzer 1997\,\cite{megretski1997iqc}）。
\end{theorem}

IQC 的力量在模块化——各经典判据都是 $\boldsymbol\Pi$ 的特例（\cref{tab:iqc}）。

\begin{table}[H]\centering\small
\caption{IQC 统一各经典稳定判据}
\label{tab:iqc}
\begin{tabular}{lll}
\toprule
$\boldsymbol\Delta$ 性质 & 对应 $\boldsymbol\Pi$ & 经典名称 \\
\midrule
有界增益 & $\mathrm{diag}(\mathbf I,-\gamma^2\mathbf I)$ & 小增益定理 \\
无源 & $\big[\begin{smallmatrix}0&\mathbf I\\\mathbf I&0\end{smallmatrix}\big]$ & Passivity 定理 \\
扇形有界 & 相应 $\boldsymbol\Pi$ & Circle criterion \\
单调递增 & Zames--Falb 乘子 & Popov/圆判据推广 \\
\bottomrule
\end{tabular}
\end{table}

IQC 严格推广 $\mu$（线性时变 $\boldsymbol\Delta$ 下退化为 $\mu$ 上界）。近年用于验证神经网络控制器鲁棒性：把 ReLU 看作满足扇形 $0\le\phi(x)\le x$、斜率 $\in[0,1]$ 的非线性算子，编码为 IQC，把“线性动力学 + 神经网络控制器”写成 $\mathbf G+\boldsymbol\Delta$ 反馈，最终归结为一个 SDP 可行性问题——为 RL 训练的策略提供部署前的安全证书。神经网络验证与鲁棒 MDP 的正式理论归 \cref{ch:clf-cbf} 与 RL 部，此处仅作统一语言的指针。KYP 引理把频域不等式 \cref{thm:iqc}(3) 与时域 LMI、无源性三者等价起来，其无源性/Port-Hamiltonian 主线归 \cref{ch:lyapunov-stability} 与 \cref{ch:operational-space}。

% ════════════════════════════════════════════════════════════════
\section{鲁棒控制 II：H∞ 综合}
\label{sec:hinf-synth}

\paragraph{与 C 章的分工。}
\cref{ch:lqr-lqg-riccati} 已把 $H_\infty$ 作为 LQR 的 minimax 推广，给出\emph{状态反馈博弈 Riccati 概览}。本节承接并深化到\emph{输出反馈 DGKF 两-Riccati 解 + 混合灵敏度 + 回路成形}——即 C 章给“$H_\infty$ 是什么”，本章给“$H_\infty$ 怎么综合 + 频域工程”。

\subsection{零和博弈解释}
LQR 等价于 $H_2$ 最优——假设扰动是白噪声、在平均意义下最优。但真实扰动往往是有结构的对手（地面突陷、负载突变）。$H_\infty$ 不做噪声统计假设，对\emph{最坏情况}优化。考虑广义被控对象 $\mathbf P$、控制器 $\mathbf K$ 构成闭环 $\mathbf T_{zw}=\mathcal F_l(\mathbf P,\mathbf K)$，$H_\infty$ 最优控制求
\begin{equation}
  \gamma^\star=\inf_{\mathbf K}\|\mathcal F_l(\mathbf P,\mathbf K)\|_\infty=\inf_{\mathbf K}\sup_\omega\bar\sigma(\mathbf T_{zw}(j\omega)).
\end{equation}

\begin{insight}
博弈解释（Başar--Bernhard\,\cite{basar1995dynamic}）：$\|\mathbf T_{zw}\|_\infty<\gamma$ 等价于零和微分博弈 $\min_u\max_w\int_0^\infty(\|z\|^2-\gamma^2\|w\|^2)\,dt<0$ 有鞍点。控制器（minimizer）选 $u$ 让性能输出 $z$ 能量最小、自然界（maximizer）选 $w$ 让 $z$ 能量最大，$\gamma$ 是扰动衰减保证“无论对手怎么干，$\|z\|/\|w\|<\gamma$”。$H_\infty$ 像下围棋（假设对手总走最坏一步），$H_2$/LQR 像和随机 AI 下棋（对手随机、只需平均表现好）。
\end{insight}

状态反馈下 Nash 均衡策略 $u^\star=-\mathbf R^{-1}\mathbf B_2^\top\mathbf P_\infty\mathbf x$、$w^\star=\gamma^{-2}\mathbf B_1^\top\mathbf P_\infty\mathbf x$，$\mathbf P_\infty$ 满足博弈 Riccati（详见 \cref{ch:lqr-lqg-riccati}）
\begin{equation}
  \mathbf A^\top\mathbf P+\mathbf P\mathbf A+\mathbf P(\gamma^{-2}\mathbf B_1\mathbf B_1^\top-\mathbf B_2\mathbf R^{-1}\mathbf B_2^\top)\mathbf P+\mathbf C_1^\top\mathbf C_1=0.
  \label{eq:game-riccati}
\end{equation}
与 LQR Riccati 的关键差别是 $+\gamma^{-2}\mathbf B_1\mathbf B_1^\top\mathbf P$ 这一\emph{正项}（代表对手破坏力），使 $\mathbf P_{H_\infty}>\mathbf P_{\text{LQR}}$、增益更激进；$\gamma\to\infty$ 时正项消失、$H_\infty\to H_2$（LQR），$\gamma\to\gamma^\star$ 时 Riccati 解爆破。实际计算 $\gamma^\star$ 用二分搜索。

\subsection{混合灵敏度}
给定权函数 $W_1,W_2,W_3$，求 $\mathbf K$ 使
\begin{equation}
  \left\|\begin{matrix}W_1\mathbf S\\W_2\mathbf K\mathbf S\\W_3\mathbf T\end{matrix}\right\|_\infty<\gamma.
  \label{eq:mixed-sens}
\end{equation}
权函数设计指导：$W_1$ 低通（低频大）约束性能（低频跟踪精度）；$W_2$ 高通或常数约束控制信号（防执行器饱和）；$W_3$ 高通（匹配不确定性权重）约束鲁棒性。这把 \cref{ch:control} 的 SISO 频域规格升级为 MIMO 的范数综合。

\subsection{DGKF 两-Riccati 解}
\begin{theorem}[DGKF 输出反馈解]
\label{thm:dgkf}
存在控制器 $\mathbf K$ 使 $\|\mathbf T_{zw}\|_\infty<\gamma$ 当且仅当（Doyle--Glover--Khargonekar--Francis 1989\,\cite{doyle1989dgkf}）：(1) 存在 $\mathbf X_\infty\ge0$ 满足控制 Riccati 且 $\mathbf A+(\gamma^{-2}\mathbf B_1\mathbf B_1^\top-\mathbf B_2\mathbf B_2^\top)\mathbf X_\infty$ 稳定；(2) 存在 $\mathbf Y_\infty\ge0$ 满足滤波 Riccati 且 $\mathbf A+(\gamma^{-2}\mathbf B_1\mathbf B_1^\top-\mathbf C_2^\top\mathbf C_2)\mathbf Y_\infty$ 稳定；(3) 谱半径耦合条件 $\rho(\mathbf X_\infty\mathbf Y_\infty)<\gamma^2$。
\end{theorem}
中心控制器具有 \textbf{LQG 分离结构}——一个状态估计器（类 Kalman）加一个状态反馈（类 LQR），但两者都被 $\gamma$ “扭曲”以提供鲁棒保证。$\gamma\to\infty$ 时 \cref{thm:dgkf} 退化为标准 LQG（\cref{ch:optimal-control-basics}）。

\subsection{Glover--McFarlane NCF 回路成形}
工业界最常用的 $H_\infty$ 流程——结合工程师的 loop-shaping 直觉与 $H_\infty$ 理论保证。对归一化互质因子 $\mathbf G=\mathbf M^{-1}\mathbf N$，闭式最大鲁棒稳定余度
\begin{equation}
  \varepsilon_{\max}=\big(1-\|[\mathbf N,\mathbf M]\|_H^2\big)^{1/2}
  \label{eq:ncf-eps}
\end{equation}
（$\|\cdot\|_H$ 为 Hankel 范数，Glover--McFarlane 1989\,\cite{glover1989ncf}），无须 $\gamma$ 迭代直接给最鲁棒控制器。三步流程：(1) \textbf{预整形} 选 $W_1,W_2$ 使 $\mathbf G_s=W_2\mathbf G W_1$ 的开环 Bode 图“长得像想要的”（$W_1$ 输入端 PI 型低通提升低频，$W_2$ 输出端 roll-off 高通衰减高频）；(2) \textbf{NCF 鲁棒稳定} 对 $\mathbf G_s$ 求最大余度控制器 $\mathbf K_\infty$（控制器阶数 = 对象阶数，不增加）；(3) \textbf{拼回} $\mathbf K=W_1\mathbf K_\infty W_2$。与混合灵敏度对照见 \cref{tab:ncf-vs-mixed}。

\begin{table}[H]\centering\small
\caption{混合灵敏度 vs Glover--McFarlane}
\label{tab:ncf-vs-mixed}
\begin{tabular}{lll}
\toprule
特征 & 混合灵敏度 & Glover--McFarlane \\
\midrule
设计自由度 & 3 个权 $W_1,W_2,W_3$ & 2 个整形权 $W_1,W_2$ \\
控制器阶数 & 对象+权之和（高） & 对象阶数（低） \\
鲁棒性保证 & 加权性能 & 互质因子不确定性 \\
工程直觉 & 需理解 $W_i$ 含义 & 直接整形 Bode 图 \\
\bottomrule
\end{tabular}
\end{table}

\textbf{ν-gap 度量}（Vinnicombe\,\cite{vinnicombe2001uncertainty}）：用 $\delta_\nu(\mathbf P_1,\mathbf P_2)$ 度量两对象在闭环意义下的距离，与回路成形结合给出鲁棒稳定的“连续性”——若控制器 $\mathbf K$ 对 $\mathbf P_1$ 的稳定余度为 $b_{\mathbf P_1,\mathbf K}$，则 $b_{\mathbf P_2,\mathbf K}\ge b_{\mathbf P_1,\mathbf K}-\delta_\nu(\mathbf P_1,\mathbf P_2)$。这把辨识误差（用 ν-gap 量）直接接到鲁棒裕度损失上。

% ════════════════════════════════════════════════════════════════
\section{鲁棒控制 III：滑模变结构控制}
\label{sec:robust-smc}

\paragraph{从频域综合回到时域非线性鲁棒。}
\cref{sec:robust-mu,sec:hinf-synth} 的 $\mu$/$H_\infty$ 是\emph{频域、线性、最坏情况}的鲁棒综合：把不确定性写成 $\boldsymbol\Delta$ 块，求一个使闭环范数最小的\emph{线性}控制器。但它有两个使用门槛——综合出的控制器阶数高、且对\emph{大范围非线性}不确定性（参数大幅变化、未建模动态、外扰）的保证依赖不确定集事先给定。本节给出一条互补路线：\textbf{滑模变结构控制}（sliding mode / variable structure control, SMC，Utkin 1977\,\cite{utkin1977variable}；Slotine--Li\,\cite{slotine1991applied}）——一个\emph{时域、非线性、对匹配不确定性结构鲁棒}的设计。它的核心思想极其朴素：不去精确建模不确定性，而是设计一个不连续的切换律，把系统状态\emph{强行约束}到一张事先选好的稳定子流形上，一旦到达，闭环动态就只由这张面的几何决定，与不确定性的具体形式无关。本节与 \cref{ch:lyapunov-stability}（Lyapunov/LaSalle）和 \cref{sec:ctrl-robot-ctc}（计算力矩）直接咬合：SMC 的到达性证明是 Lyapunov 法的有限时间变体，而其机器人实现正是计算力矩的鲁棒化。

\subsection{滑模面与到达条件}
考虑单输入非线性系统（多输入逐通道推广）
\begin{equation}
  x^{(n)}=f(\mathbf x)+b(\mathbf x)\,u+d(t),\qquad \mathbf x=[x,\dot x,\dots,x^{(n-1)}]^\top,
  \label{eq:smc-plant}
\end{equation}
其中 $f,b$ 名义已知但有不确定性，$d(t)$ 为外扰，$b(\mathbf x)\ge b_0>0$。给定参考 $x_d(t)$，定义跟踪误差 $e=x-x_d$。\textbf{滑模面}取误差及其导数的 Hurwitz 线性组合
\begin{equation}
  s(\mathbf x,t)=\Big(\frac{d}{dt}+\lambda\Big)^{n-1}e
  =e^{(n-1)}+\binom{n-1}{1}\lambda\,e^{(n-2)}+\cdots+\lambda^{n-1}e,\qquad \lambda>0.
  \label{eq:smc-surface}
\end{equation}
关键洞察：方程 $s(\mathbf x,t)=0$ 是状态空间里一张 $(n-1)$ 维曲面，而把 $s=0$ 看成关于 $e$ 的微分方程，其特征根全在 $-\lambda$（$n-1$ 重），故\emph{只要状态停留在 $s=0$ 上，误差 $e$ 必按时间常数 $1/\lambda$ 指数收敛到零}。于是 $n$ 阶跟踪问题被降维成两个子问题：(i) 把 $s$ 驱到 $0$（\emph{到达阶段}）；(ii) 维持在 $s=0$（\emph{滑动阶段}），后者的收敛由面的设计自动保证。

\begin{definition}[到达条件]
\label{def:smc-reaching}
若控制律使 $s$ 满足
\begin{equation}
  s\,\dot s<0\quad(\forall\, s\neq0),\qquad\text{更强的}\quad \tfrac12\frac{d}{dt}s^2=s\dot s\le-\eta|s|,\ \ \eta>0,
  \label{eq:smc-reaching}
\end{equation}
则称 $s=0$ 满足\textbf{到达条件}（$\eta$-reachability）。前者保证 $s$ 单调趋零，后者（$\eta$-可达，constant-reaching）进一步保证\emph{有限时间}到达。
\end{definition}

\begin{theorem}[有限时间到达]
\label{thm:smc-finite-reach}
若 $\tfrac12\frac{d}{dt}s^2\le-\eta|s|$，则从任意初值 $s(0)$ 出发，$s$ 在
\begin{equation}
  t_r\le \frac{|s(0)|}{\eta}
  \label{eq:smc-treach}
\end{equation}
内到达 $s=0$，此后永远停留（滑模不变性）。
\end{theorem}
\begin{proof}[骨架]
取 $V=\tfrac12 s^2\ge0$，则 $\dot V=s\dot s\le-\eta|s|=-\eta\sqrt{2V}$。对 $s>0$ 区域 $\frac{d}{dt}|s|=\mathrm{sign}(s)\dot s\le-\eta$，积分得 $|s(t)|\le|s(0)|-\eta t$，故 $|s|$ 在 $t_r\le|s(0)|/\eta$ 触零。注意这\emph{快于}指数收敛——指数律 $\dot V=-2\eta V$ 只渐近趋零，而 $\dot V\le-\eta\sqrt{2V}$ 的 $\sqrt{V}$ 右端使 $V$ 有限时间归零，这正是不连续切换项相对连续反馈的本质优势。完整有限时间稳定性理论见 \cref{ch:lyapunov-stability}。
\end{proof}

\subsection{等效控制与切换控制}
在滑模面上 $\dot s=0$，由此解出维持滑动所需的\textbf{等效控制} $u_{\mathrm{eq}}$（Utkin 的等效控制法）。以二阶 $n=2$、$s=\dot e+\lambda e$ 为例，$\dot s=\ddot e+\lambda\dot e=\ddot x-\ddot x_d+\lambda\dot e=f+bu+d-\ddot x_d+\lambda\dot e$，令名义部分 $\dot s=0$（设 $d=0$、$f,b$ 取名义 $\hat f,\hat b$）得
\begin{equation}
  u_{\mathrm{eq}}=\hat b^{-1}\big(-\hat f+\ddot x_d-\lambda\dot e\big).
  \label{eq:smc-ueq}
\end{equation}
$u_{\mathrm{eq}}$ 是把状态\emph{保持}在面上的“连续等效输入”，几何上等于真实不连续控制在切换面上的局部平均（Filippov 意义，见 \cref{ch:lyapunov-stability} 对 $\mathrm{sign}$ 破坏 Lipschitz 的讨论）。完整控制律再叠加一个\textbf{切换控制}把状态拉向面并压制不确定性
\begin{equation}
  u=u_{\mathrm{eq}}-\hat b^{-1}K\,\mathrm{sign}(s),\qquad K>0.
  \label{eq:smc-law}
\end{equation}

\begin{theorem}[匹配不确定性下的鲁棒到达]
\label{thm:smc-robust}
设不确定性\textbf{匹配}（与控制同通道）：真实 $\ddot s=-\hat b\,\hat b^{-1}K\,\mathrm{sign}(s)+\Delta(\mathbf x,t)$，其中集总扰动 $|\Delta|\le D(\mathbf x,t)$ 有已知上界。若切换增益取
\begin{equation}
  K\ge D+\eta,
  \label{eq:smc-gain}
\end{equation}
则 $\tfrac12\frac{d}{dt}s^2=s\big(-K\,\mathrm{sign}(s)+\Delta\big)\le-(K-D)|s|\le-\eta|s|$，到达条件成立，闭环对所有满足 $|\Delta|\le D$ 的不确定性\emph{鲁棒}。
\end{theorem}

\begin{insight}
SMC 鲁棒性的本质是\textbf{结构性}而非定量调参：只要不确定性\emph{匹配}且其界 $D$ 已知，切换增益 $K\ge D+\eta$ 就把它“吃掉”，滑动阶段的闭环动态\emph{完全不依赖} $f,b$ 的真实值——这叫\textbf{滑模不变性}（invariance, 强于鲁棒性的扰动衰减）。这与 \cref{sec:hinf-synth} 的 $H_\infty$ 形成鲜明对照：$H_\infty$ 把扰动\emph{衰减}到 $\gamma$ 倍（$\|z\|<\gamma\|w\|$，仍残留），SMC 对匹配扰动做到\emph{完全抵消}（$s\equiv0$ 上 $e$ 的动态与扰动无关）。代价是这只对\emph{匹配}不确定性成立：不匹配（unmatched）部分 SMC 无法直接抵消，需配合积分滑模或 backstepping。
\end{insight}

\subsection{抖振及其治理}
\cref{eq:smc-law} 的 $\mathrm{sign}(s)$ 在 $s=0$ 处不连续，理想上以无穷频率切换。真实系统有采样延迟、执行器惯性与未建模高频动态，使轨迹在面两侧\emph{高频来回穿越}而非完美滑动——这就是\textbf{抖振}（chattering）。抖振激励柔性模态、磨损执行器、产生噪声，是 SMC 工程落地的头号障碍。三类治理手段：

\paragraph{(1) 边界层 / 饱和函数。}
在面附近开一条厚度 $\phi$ 的边界层，用连续的饱和 $\mathrm{sat}(s/\phi)$ 替换 $\mathrm{sign}(s)$
\begin{equation}
  \mathrm{sat}(s/\phi)=
  \begin{cases}
    s/\phi, & |s|\le\phi,\\
    \mathrm{sign}(s), & |s|>\phi,
  \end{cases}
  \qquad
  u=u_{\mathrm{eq}}-\hat b^{-1}K\,\mathrm{sat}(s/\phi).
  \label{eq:smc-sat}
\end{equation}
边界层内控制律退化为\emph{高增益线性反馈}（斜率 $K/\phi$），轨迹不再切换而是被吸入薄层。代价是放弃了面上的理想不变性：稳态跟踪误差不再为零，而是落在量级 $|e|\sim O(\phi/\lambda^{n-1})$ 的邻域内（边界层越厚抖振越小、精度越低，是\textbf{抖振-精度权衡}）。$\phi$ 应按可容忍的跟踪误差与未建模动态带宽折中选取（Slotine--Li\,\cite{slotine1991applied}：$\phi$ 与跟踪精度成正比，可用时变 $\phi(t)$ 随增益自适应）。

\paragraph{(2) 高阶滑模。}
边界层以牺牲精度换平滑。\textbf{高阶滑模}（higher-order SMC，Levant\,\cite{levant1993sliding}）另辟蹊径：把不连续性\emph{推到控制的高阶导数上}，使实际控制信号连续。$r$ 阶滑模不仅令 $s=0$，还令 $\dot s=\cdots=s^{(r-1)}=0$，对相对阶为 $r$ 的输出实现有限时间收敛且控制连续（不连续项藏在 $u^{(r-1)}$ 里）。

\paragraph{(3) 超螺旋（二阶滑模）。}
最常用的高阶滑模是\textbf{超螺旋算法}（super-twisting algorithm, STA，Levant 1993\,\cite{levant1993sliding}），它只需 $s$ 的测量（无须 $\dot s$），对相对阶 1 的 $s$ 给出连续控制并在有限时间内令 $s=\dot s=0$：
\begin{align}
  u&=-k_1|s|^{1/2}\,\mathrm{sign}(s)+v,\label{eq:sta-u}\\
  \dot v&=-k_2\,\mathrm{sign}(s).\label{eq:sta-v}
\end{align}
第一项是连续的（$|s|^{1/2}$ 在 $s=0$ 连续），第二项把不连续的 $\mathrm{sign}$ 放进\emph{积分器} $v$，故输出到执行器的 $u$ 连续——抖振被本质性削弱而非仅仅抹平。对满足 $|\dot\Delta|\le L$ 的扰动，增益取 $k_1>1.5\sqrt L,\ k_2>1.1\,L$（Moreno--Osorio 严格 Lyapunov 充分条件的常用工程取值；精确常数依扰动模型而定）即保证有限时间收敛。STA 的几何图像是状态在 $(s,\dot s)$ 相平面里螺旋收敛到原点（故名 super-twisting），收敛时间有限。

\begin{table}[H]\centering\small
\caption{三类抖振治理手段对照}
\label{tab:smc-chatter}
\begin{tabular}{llll}
\toprule
手段 & 控制连续性 & 滑模精度 & 代价/前提 \\
\midrule
边界层 $\mathrm{sat}(s/\phi)$ & 连续 & $O(\phi)$ 邻域（非理想） & 牺牲稳态精度 \\
高阶滑模（$r$ 阶） & 高阶导连续 & 理想有限时间 & 需高阶导数信息 \\
超螺旋 STA & 连续（$u$） & 理想有限时间 & 需扰动导数界 $|\dot\Delta|\le L$ \\
\bottomrule
\end{tabular}
\end{table}

\subsection{应用于机器人轨迹跟踪}
把 SMC 接到机器人动力学 $\mathbf M(\mathbf q)\ddot{\mathbf q}+\mathbf C(\mathbf q,\dot{\mathbf q})\dot{\mathbf q}+\mathbf g(\mathbf q)=\boldsymbol\tau$（\cref{sec:ctrl-robot-ctc}）。设期望轨迹 $\mathbf q_d(t)$，误差 $\mathbf e=\mathbf q-\mathbf q_d$，取\textbf{向量滑模面}
\begin{equation}
  \mathbf s=\dot{\mathbf e}+\boldsymbol\Lambda\mathbf e,\qquad \boldsymbol\Lambda=\boldsymbol\Lambda^\top\succ0.
  \label{eq:smc-robot-surface}
\end{equation}
这正是 \cref{eq:smc-surface} 在 $n=2$ 的多变量版，也与 \cref{ch:force-wbc} 中 Slotine--Li 自适应控制的滑模变量同源——区别在那里用自适应估计参数，这里用切换项压制参数误差。定义参考速度 $\dot{\mathbf q}_r=\dot{\mathbf q}_d-\boldsymbol\Lambda\mathbf e$，则 $\mathbf s=\dot{\mathbf q}-\dot{\mathbf q}_r$。控制律取计算力矩结构再加切换项（$\hat{\mathbf M},\hat{\mathbf C},\hat{\mathbf g}$ 为名义模型）
\begin{equation}
  \boldsymbol\tau=\hat{\mathbf M}(\mathbf q)\ddot{\mathbf q}_r+\hat{\mathbf C}(\mathbf q,\dot{\mathbf q})\dot{\mathbf q}_r+\hat{\mathbf g}(\mathbf q)-\mathbf K_s\,\mathrm{sign}(\mathbf s),\qquad \mathbf K_s=\mathrm{diag}(K_{s,i})\succ0.
  \label{eq:smc-robot-law}
\end{equation}

\begin{theorem}[机器人 SMC 跟踪误差收敛]
\label{thm:smc-robot}
取 Lyapunov 候选 $V=\tfrac12\mathbf s^\top\mathbf M(\mathbf q)\mathbf s$。利用反对称性质 $\dot{\mathbf M}-2\mathbf C$ 反对称（$\mathbf s^\top(\dot{\mathbf M}-2\mathbf C)\mathbf s=0$，见 \cref{sec:ctrl-robot-ctc}），沿闭环
\begin{equation}
  \dot V=\mathbf s^\top\mathbf M\dot{\mathbf s}+\tfrac12\mathbf s^\top\dot{\mathbf M}\mathbf s
       =\mathbf s^\top\big(\boldsymbol\tau-\mathbf M\ddot{\mathbf q}_r-\mathbf C\dot{\mathbf q}_r-\mathbf g\big)+\mathbf s^\top\mathbf C\mathbf s+\tfrac12\mathbf s^\top\dot{\mathbf M}\mathbf s
       =\mathbf s^\top\big(\boldsymbol\eta-\mathbf K_s\,\mathrm{sign}(\mathbf s)\big),
  \label{eq:smc-robot-vdot}
\end{equation}
其中 $\boldsymbol\eta=(\hat{\mathbf M}-\mathbf M)\ddot{\mathbf q}_r+(\hat{\mathbf C}-\mathbf C)\dot{\mathbf q}_r+(\hat{\mathbf g}-\mathbf g)$ 是集总建模误差。若逐分量取 $K_{s,i}\ge|\eta_i|+\eta_0$（$\eta_0>0$），则 $\dot V\le-\eta_0\sum_i|s_i|<0$，故 $\mathbf s\to\mathbf 0$ 有限时间到达；到达后由 \cref{eq:smc-robot-surface} 知 $\dot{\mathbf e}=-\boldsymbol\Lambda\mathbf e$，$\mathbf e\to\mathbf 0$ 指数收敛。
\end{theorem}

\begin{example}[二连杆机械臂的负载不确定]
末端负载未知（$0$--$5\,$kg）使 $\mathbf g(\mathbf q)$ 与 $\mathbf M(\mathbf q)$ 大幅变化。纯计算力矩（\cref{sec:ctrl-robot-ctc}）按名义负载补偿，满载时残留重力矩使稳态误差显著；SMC 用 \cref{eq:smc-robot-law}，只要切换增益覆盖最大负载引起的 $|\eta_i|$，跟踪误差就被压到边界层精度内，\emph{无须知道负载具体值}。实现时把 $\mathrm{sign}$ 换成 $\mathrm{sat}(s_i/\phi_i)$ 消抖（\cref{eq:smc-sat}），稳态误差换得平滑力矩；若要求零稳态误差则用超螺旋 \cref{eq:sta-u,eq:sta-v} 逐关节实现。
\end{example}

\begin{pitfall}[SMC 三个易错点]
\begin{enumerate}
  \item \textbf{增益盲目调大}——$K$ 越大鲁棒裕度越大，但抖振幅度与执行器应力也随 $K$ 增大；$K$ 应贴着不确定性界 $D+\eta$ 取，而非越大越好。
  \item \textbf{忽视不匹配不确定性}——\cref{thm:smc-robust} 的不变性\emph{只对匹配（与 $\mathbf u$ 同通道）不确定性}成立。柔性关节里弹性变形是不匹配的，直接 SMC 无效，须用积分滑模或对负载侧建面。
  \item \textbf{在仿真里信任不连续解}——$\mathrm{sign}(s)$ 破坏 Lipschitz，定步长积分器会产生虚假高频极限环或数值发散（见 \cref{ch:lyapunov-stability}）。验证 SMC 必须用边界层光滑化或事件检测积分器，否则“仿真抖振”可能是数值假象而非真实现象。
\end{enumerate}
\end{pitfall}

\subsection{与本章 H∞/μ 的对照}
SMC 与 \cref{sec:robust-mu,sec:hinf-synth} 的频域综合是鲁棒控制的两条互补主线，定位差异见 \cref{tab:smc-vs-hinf}。

\begin{table}[H]\centering\small
\caption{滑模 vs $H_\infty$/$\mu$：两条鲁棒路线的定位}
\label{tab:smc-vs-hinf}
\begin{tabular}{lll}
\toprule
维度 & 滑模 SMC & $H_\infty$/$\mu$ 综合 \\
\midrule
工作域 & 时域、非线性 & 频域、线性（LFT） \\
鲁棒对象 & \textbf{匹配}不确定性（界已知） & 结构化 $\boldsymbol\Delta$（范数有界） \\
保证类型 & 不变性（完全抵消） & 衰减（$\|z\|<\gamma\|w\|$） \\
控制器复杂度 & 低（不连续律 + 名义模型） & 高（阶数 = 对象+权） \\
主要代价 & 抖振、需不确定性界 & 综合保守、需频域建模 \\
适用场景 & 大范围参数变化、强外扰 & 多目标频域规格、MIMO 耦合 \\
\bottomrule
\end{tabular}
\end{table}

\begin{insight}
两者并非对立而是\textbf{可叠加}：工程上常用“$H_\infty$ 做名义回路成形保证频域性能 + SMC/超螺旋做内环压制匹配的大不确定性”。判据上也有统一视角——SMC 的到达条件 $s\dot s<0$ 是一个时域 Lyapunov 不等式，而小增益（\cref{thm:small-gain}）/$\mu$（\cref{thm:mu-rs}）是频域判据；两者经 IQC/KYP 可纳入同一“耗散性”语言（\cref{ch:lyapunov-stability}）。选择准则：不确定性\emph{可结构化建模为频域 $\boldsymbol\Delta$}且要 MIMO 多目标 $\Rightarrow$ 走 $H_\infty$/$\mu$；不确定性\emph{大但匹配且界可估}、要简单实现与强扰动抑制 $\Rightarrow$ 走 SMC。把 \cref{eq:mixed-sens} 的混合灵敏度与 \cref{thm:dgkf} 的 DGKF 中心控制器与本节的 \cref{eq:smc-robot-law} 并置，就看清了“频域衰减”与“时域不变”是同一鲁棒目标的两种数学化。
\end{insight}

% ════════════════════════════════════════════════════════════════
\section{频域多变量分析}
\label{sec:mimo-freq}

\paragraph{SISO 工具的 MIMO 升级。}
SISO 的 Nyquist/Bode 已在 \cref{ch:control}，此处只补 MIMO 推广。本质困难：SISO 的 $|G(j\omega)|$ 唯一确定增益，但 MIMO 的增益取决于\emph{输入方向}——2 输入 2 输出系统沿不同输入方向增益可差 100 倍。

\subsection{奇异值与方向性}
在每个频率对 $\mathbf G(j\omega)$ 做 SVD $\mathbf G(j\omega)=\mathbf U(\omega)\boldsymbol\Sigma(\omega)\mathbf V(\omega)^*$：最大奇异值 $\bar\sigma=\max_{\|u\|=1}\|\mathbf Gu\|$ 是最坏方向增益、最小奇异值 $\underline\sigma$ 是最好方向增益下界、条件数 $\bar\sigma/\underline\sigma$ 度量方向病态（机械臂近奇异位形时雅可比条件数趋无穷，沿某些方向几乎不能运动）。\textbf{回路整形准则}（$\mathbf L=\mathbf G\mathbf K$）：低频要 $\underline\sigma(\mathbf L)$ 大（所有方向都有足够增益跟踪/抑扰）、高频要 $\bar\sigma(\mathbf L)$ 小（最坏方向不放大噪声）、中频两者冲突是设计核心挑战。只看 $\bar\sigma$ 不看 $\underline\sigma$ 会设计出某方向跟踪完美、另一方向完全失控的控制器。

\subsection{RGA 与输入输出配对}
对方阵稳态增益 $\mathbf G(0)$，相对增益阵列（Bristol 1966\,\cite{bristol1966rga}）
\begin{equation}
  \boldsymbol\Lambda=\mathbf G(0)\circ(\mathbf G(0)^{-1})^\top
\end{equation}
（$\circ$ 为 Hadamard 积），$\lambda_{ij}$ 衡量“输入 $j$ 对输出 $i$ 的相对影响”（计入其他回路交互）：$\lambda_{ij}=1$ 无交互（单回路即可）、$\gg1$ 正交互且敏感、$<0$ 负交互（其他回路打开时增益反向）、$\approx0$ 无直接通路。\textbf{配对规则}：选 $\lambda_{ij}$ 接近 1 的对配对，避开负或远大于 1 者。双关节机械臂末端 $(x,y)$ 控制：若关节 1 主要影响 $x$、关节 2 主要影响 $y$，RGA 接近单位阵、对角控制器有效；若近奇异，RGA 远偏单位阵、须耦合控制器。

\subsection{水床效应（MIMO/RHP 极点版）}
MIMO 灵敏度 $\mathbf S=(\mathbf I+\mathbf L)^{-1}$、补灵敏度 $\mathbf T=\mathbf L(\mathbf I+\mathbf L)^{-1}$，$\mathbf S+\mathbf T=\mathbf I$（\cref{eq:ctrl-S-plus-T}）仍成立但 $\mathbf S,\mathbf T$ 不对易。频域规格表见 \cref{tab:freq-spec}。

\begin{table}[H]\centering\small
\caption{MIMO 频域控制规格}
\label{tab:freq-spec}
\begin{tabular}{lll}
\toprule
规格 & 数学表达 & 含义 \\
\midrule
性能（低频） & $\|W_1\mathbf S\|_\infty<1$ & 加权灵敏度有界 \\
控制努力 & $\|W_2\mathbf K\mathbf S\|_\infty<1$ & 控制信号不过大 \\
鲁棒性（高频） & $\|W_3\mathbf T\|_\infty<1$ & 不确定性不破坏稳定 \\
\bottomrule
\end{tabular}
\end{table}

频域控制设计最深刻的基本限制是\textbf{Bode 灵敏度积分}（水床效应）：

\begin{theorem}[Bode 灵敏度积分]
\label{thm:bode-integral}
对 SISO 稳定闭环，若开环 $L(s)$ 的相对阶 $\ge2$，则（Freudenberg--Looze 1985\,\cite{freudenberg1985rhp}；原始思想 Bode 1945\,\cite{bode1945network}）
\begin{equation}
  \int_0^\infty\ln|S(j\omega)|\,d\omega=\pi\sum_i\mathrm{Re}(p_i),
  \label{eq:bode-integral}
\end{equation}
其中 $p_i$ 为 $L(s)$ 的 RHP 极点。
\end{theorem}

$\ln|S|$ 的积分是守恒量——像按压水床，一处压下（$|S|<1$，抑扰）另一处必鼓起（$|S|>1$，放大）。三种情形：稳定（无 RHP 极点）积分 $=0$（零和）；不稳定（有 RHP 极点）$>0$（放大面积严格大于抑制面积，稳定化有代价）；非最小相位（RHP 零点 $z$）有额外 Poisson 加权约束 $\int_0^\infty\ln|S(j\omega)|\tfrac{2\mathrm{Re}(z)}{|j\omega-z|^2}d\omega=\pi\ln\big|\tfrac{\bar z+p}{z-p}\big|$，Poisson 核在 $\omega=\mathrm{Im}(z)$ 附近集中，意味着 RHP 零点对应频率附近 $|S|$ 必须 $>1$。

\begin{proof}[水床积分证明概要]
$S(s)=1/(1+L(s))$，闭环稳定时 $S$ 在 RHP 解析且非零，故 $\ln S$ 也在 RHP 解析。对 Nyquist D-围道用 Cauchy 积分：虚轴部分贡献 $\int_{-\infty}^\infty\ln S(j\omega)\,d\omega$ 的相关项，大半圆部分因相对阶 $\ge2$ 使 $S(s)\to1$、$\ln S\to0$ 故贡献为零。\emph{关键（须正确表述）}：开环 RHP 极点 $p_i$ 是 $L$ 的极点，故在 $p_i$ 处 $L\to\infty$、$S(p_i)=1/(1+L(p_i))\to0$，即 $p_i$ 是 $S$ 的\textbf{零点}（不是闭环极点条件 $1+L=0$）。$\ln S$ 在零点 $p_i$ 处的留数给出 $\pi\,\mathrm{Re}(p_i)$ 的贡献，对所有 RHP 极点求和即得 \cref{eq:bode-integral}。完整细节见 Skogestad--Postlethwaite\,\cite{skogestad2005multivariable} Ch.5.3 或 Doyle--Francis--Tannenbaum\,\cite{doyle1990feedback} Ch.6。
\end{proof}

\begin{insight}
水床效应给出机器人设计的\textbf{物理硬限制}：RHP 零点 $z\Rightarrow$ 闭环带宽上限 $\omega_c\lesssim z/2$；RHP 极点 $p\Rightarrow$ 带宽下限 $\omega_c\gtrsim2p$；若 $z<4p$ 可能无法同时满足——系统结构上不可镇定到满意性能。倒立摆 RHP 极点 $p\approx\sqrt{g/l}$，$l=1\,$m 的腿 $p\approx3.1\,$rad/s，带宽至少 $2p\approx6.3\,$rad/s，留 10 倍裕度 $\Rightarrow$ MPC 更新率至少约 10\,Hz。这不是工程选择，是物理必然。好的控制设计不是消除水床，而是把鼓起放在不重要的频段（高频噪声区）。
\end{insight}

\subsection{广义 Nyquist 与特征轨迹}
\textbf{广义 Nyquist 定理}：MIMO 闭环稳定 $\iff\det(\mathbf I+\mathbf L(s))$ 沿 Nyquist D-围道绕\emph{原点}的圈数 $N=-P$（$P$ = 开环 RHP 极点数）。\textbf{特征轨迹法}（MacFarlane--Postlethwaite 1977\,\cite{macfarlane1977characteristic}）：算 $\mathbf L(j\omega)$ 的特征值 $\lambda_i(\mathbf L(j\omega))$ 随 $\omega$ 变化，绘 $n$ 条 characteristic loci，所有轨迹绕 $-1$ 的逆时针总圈数 $=P$ 则稳定。\textbf{方向 Nyquist}：在选定方向 $\mathbf v$ 上考察 $\mathbf v^*\mathbf L(j\omega)\mathbf v$，揭示特定控制方向的稳定裕度——对 pitch/roll 解耦分析尤其有用。Nichols 图（横轴相位、纵轴幅度）上恒定 $|S|$、$|T|$ 的等值线（M/N 圆）成标准曲线，控制设计即“让 Nichols 曲线避开 $|S|>S_{\max}$ 区域”，是 QFT（Horowitz\,\cite{horowitz1963qft}）的核心工具。

% ════════════════════════════════════════════════════════════════
\section{桥接与小结}
\label{sec:robust-summary}

本章从“模型从哪来”出发，沿辨识 $\to$ 数据驱动控制 $\to$ 鲁棒综合 $\to$ 频域验证一条主线，把 \cref{ch:optimal-control-basics} 的 LQR/LQG 的三个缺口逐一补完。回扣开篇：Doyle 1978\,\cite{doyle1978guaranteed} 的 LQG 反例正是本章存在的理由——$H_2$ 最优不蕴含鲁棒，必须有 $\mu$/$H_\infty$/IQC 这一层；DGKF 的中心控制器（\cref{thm:dgkf}）则把 LQG 的分离结构在 $\gamma$ 扭曲下重新拾回，给出了带鲁棒保证的输出反馈。

\begin{insight}[鲁棒控制 $\leftrightarrow$ RL 鲁棒性的启发式对应]
在 sim-to-real 中，参数不确定性（摩擦、质量、延迟）对应结构化 $\boldsymbol\Delta$，domain randomization 范围\emph{启发式地}对应不确定集半径、与 $H_\infty$ 的 $1/\gamma$ 定性呼应：范围太大（相当于 $\gamma$ 太小）则策略过保守甚至不收敛（Riccati 无解），太小则不够鲁棒。须强调这是\emph{类比而非定理}——domain randomization 是鲁棒 MDP（Iyengar 2005\,\cite{iyengar2005robustdp}；Nilim--El Ghaoui 2005\,\cite{nilim2005robust}）的采样近似（soft-robust），与 $\gamma$ 只是定性对应，不应写成等号。鲁棒 MDP 的 robust Bellman 收缩、神经网络验证的 SDP 证书等正式理论让位 \cref{ch:clf-cbf} 与 RL 部。
\end{insight}

\begin{pitfall}[常见误解]
\begin{enumerate}
  \item \textbf{“数据越多辨识越准”}——错。若轨迹不激励某些参数方向，再多数据也辨不出；关键是信息矩阵条件数（\cref{sec:robot-id}）。
  \item \textbf{“辨识精度越高控制越好”}——错。边际收益递减，精度应匹配控制需求（鲁棒控制器可容忍 $10\%$--$20\%$ 误差）。
  \item \textbf{“RLS 与 Kalman 无关”}——错，RLS 就是 Kalman 的特例（\cref{thm:rls-kalman}）。
  \item \textbf{“小增益对结构不确定性也是充要的”}——错，满块才充要，结构块下小增益保守，须用 $\mu$（\cref{thm:mu-rs}）。
  \item \textbf{“D-scaling 上界出自 Doyle 1982”}——错，$\mu$ 定义归 Doyle 1982，D-scaling 上界理论归 Packard--Doyle 1993。
  \item \textbf{“可以让所有频段 $|S|<1$”}——错，水床积分守恒（\cref{thm:bode-integral}），先问“愿在哪个频段牺牲性能”。
  \item \textbf{“domain randomization 范围 = $1/\gamma$ 是等式”}——错，仅是启发式对应，DR 是鲁棒 MDP 的采样近似。
\end{enumerate}
\end{pitfall}

\subsection*{延伸阅读}
辨识：Ljung\,\cite{ljung1999sysid}（PEM/模型族/渐近的标准教材）、Söderström--Stoica\,\cite{soderstrom1989sysid}、机器人动力学辨识 Atkeson--An--Hollerbach\,\cite{atkeson1986inertial}/Gautier--Khalil\,\cite{gautier1991identifiable}/Swevers\,\cite{swevers1997optimal}、频域辨识 Pintelon--Schoukens\,\cite{pintelon2012freqdomain}。数据驱动控制：Willems 等\,\cite{willems2005persistency}、DeePC\,\cite{coulson2019deepc,coulson2019regularized}、数据信息性\,\cite{vanwaarde2020informativity}、De Persis--Tesi\,\cite{depersis2020datadriven}、带稳定保证的数据驱动 MPC\,\cite{berberich2021ddmpc}、低秩逼近\,\cite{markovsky2019lowrank}、SLS\,\cite{anderson2019sls,dean2020sample}。鲁棒/H∞/μ：圣经 Zhou--Doyle--Glover\,\cite{zhou1996robust}、凸方法 Dullerud--Paganini\,\cite{dullerud2000convex}、$\mu$ 定义\,\cite{doyle1982mu} 与 D-scaling\,\cite{packard1993mu}、IQC\,\cite{megretski1997iqc}、神经网络验证\,\cite{fazlyab2019nn}、DGKF\,\cite{doyle1989dgkf}、NCF 回路成形\,\cite{glover1989ncf,mcfarlane1990loopshaping}、H∞ 博弈\,\cite{basar1995dynamic}、ν-gap\,\cite{vinnicombe2001uncertainty}。频域/MIMO：工程圣经 Skogestad--Postlethwaite\,\cite{skogestad2005multivariable}、H∞ SISO 入门 Doyle--Francis--Tannenbaum\,\cite{doyle1990feedback}、RGA\,\cite{bristol1966rga}、特征轨迹\,\cite{macfarlane1977characteristic}、QFT\,\cite{horowitz1963qft}、LQR 裕度\,\cite{kalman1964optimal} 与 LQG 反例\,\cite{doyle1978guaranteed}、柔性关节\,\cite{spong1987flexible}。

\noindent\textbf{启下.}\;有了辨识给出的模型与鲁棒/频域的验证工具，下一步是把它们用于具体的安全与最优控制综合：\cref{ch:clf-cbf} 用 CLF/CBF 把安全约束写成 QP（水床/带宽硬限制与安全滤波互引）、\cref{ch:mpc-advanced} 把 DeePC 与鲁棒 MPC（Tube）展开、\cref{ch:hj-reachability} 从可达集角度给出与鲁棒稳定不同的安全路线。辨识的回归矩阵 $\mathbf Y$ 与柔性关节模型则在机械臂动力学（\cref{ch:operational-space}）与力位混合/WBC（\cref{ch:force-wbc}）中复用。
